\documentclass[aps,preprint]{revtex4}%
\usepackage{amsfonts}
\usepackage{amsmath}
\usepackage{amssymb}
\usepackage{graphicx}%
\setcounter{MaxMatrixCols}{30}
%TCIDATA{OutputFilter=latex2.dll}
%TCIDATA{Version=4.10.0.2363}
%TCIDATA{CSTFile=revtex4.cst}
%TCIDATA{Created=Friday, August 29, 2003 10:08:38}
%TCIDATA{LastRevised=Saturday, July 24, 2004 09:27:28}
%TCIDATA{<META NAME="GraphicsSave" CONTENT="32">}
%TCIDATA{<META NAME="DocumentShell" CONTENT="Articles\SW\REVTeX 4">}
%TCIDATA{Language=American English}
\newtheorem{theorem}{Theorem}
\newtheorem{acknowledgement}[theorem]{Acknowledgement}
\newtheorem{algorithm}[theorem]{Algorithm}
\newtheorem{axiom}[theorem]{Axiom}
\newtheorem{claim}[theorem]{Claim}
\newtheorem{conclusion}[theorem]{Conclusion}
\newtheorem{condition}[theorem]{Condition}
\newtheorem{conjecture}[theorem]{Conjecture}
\newtheorem{corollary}[theorem]{Corollary}
\newtheorem{criterion}[theorem]{Criterion}
\newtheorem{definition}[theorem]{Definition}
\newtheorem{example}[theorem]{Example}
\newtheorem{exercise}[theorem]{Exercise}
\newtheorem{lemma}[theorem]{Lemma}
\newtheorem{notation}[theorem]{Notation}
\newtheorem{problem}[theorem]{Problem}
\newtheorem{proposition}[theorem]{Proposition}
\newtheorem{remark}[theorem]{Remark}
\newtheorem{solution}[theorem]{Solution}
\newtheorem{summary}[theorem]{Summary}
\newenvironment{proof}[1][Proof]{\noindent\textbf{#1.} }{\ \rule{0.5em}{0.5em}}
\begin{document}
\preprint{ }
\title{Correlated One Particle Wave Functions}
\author{B. Weiner}
\affiliation{Department of Physics, Pennsylvania State University, DuBois PA 15801}
\author{J. V. Ortiz}
\affiliation{Department of Chemistry, Kansas State University, Manhattan, KS 66506-3701}
\pacs{}

\begin{abstract}


\end{abstract}
\date{06/29/2004}
\startpage{1}
\maketitle
\tableofcontents


\section{ Introduction \label{intro}}

Independent Particle States (IPS's) have proved a mainstay of molecular
multi-particle quantum mechanics both in physical modelling and numerical
applications for many decades. The major attraction of these states lies in
the conceptually concrete and chemically understandable pictures of molecular
structure and properties that they provide. Apart from an initial choice of
$s$ atomic orbitals and possibly a decision about the molecular point group
and spin symmetry, one can obtain an unbiased description of the electronic
states of a molecule by finding an IPS\ that minimizes the energy via the Self
Consistent Field (SCF)\ equations. These equations provide a set of one
electron states $\left\{  \left.  \left\vert \varphi_{j}\right\rangle
\right\vert 1\leq j\leq2s\right\}  $, that are eigenstates of a self
consistent Fock operator corresponding to the eigenvalues $\left\{  \left.
\varepsilon_{j}\right\vert 1\leq j\leq2s\right\}  $ and are unique up to
unitary transformations that mix degenerate eigenstates. If the IPS\ describes
$M$ electrons then in general it is determined by the one-electron eigenstates
corresponding to the $M$ lowest eigenvalues $\left\{  \left.  \varepsilon
_{j}\right\vert 1\leq j\leq M\right\}  $. The wave functions of each of these
one electron states can then be analyzed in terms of a basis of atomic
orbitals, producing the familiar picture of one electron Probability
Distributions (PD's) labelled by one electron energies, that has been
invaluable in the understanding of chemical processes. The physical meaning of
the energies $\left\{  \left.  \varepsilon_{j}\right\vert 1\leq j\leq
M\right\}  $ can be deduced by considering the electron propagator
\cite{ortiz}, which shows that they correspond to ionization potentials.

It is often the case, however, that the IPS\ model is not sufficiently
accurate to predict or explain the finer details of chemical reaction
mechanisms or spectroscopic observations. In order to achieve this one must
take into account electron-electron correlation. Traditionally this is
accomplished by perturbational or configurational interaction methods that
increase numerical accuracy at the cost of a clear, tangible, chemical model
of the molecular system. The clarity is lost as pictorial models need be based
on one particle properties, but in general one particle properties do not
determine correlated multi-electron states. In particular one cannot use
chemical intuition to rearrange the positions of atoms and bonds and be
certain that the new conformation actually corresponds to a multi-electron
state in contrast to the case of states determined by one particle properties.

However, we have described \cite{Weiner & Ortiz 2003} a class of correlated
states that generalize the IPS model, explicitly contain a description of
correlation effects, and are the \emph{most general} states that produce a
\emph{one particle} theory of electrons i.e. multi-particle states determined
by one particle properties. These are the Antisymmetrized Geminal Power (AGP)
states for an even number, $2N$, of electrons and the Generalized
Antisymmetrized Geminal Power (GAGP) ones for an odd number, $2N+1$. In the
IPS\ case the $2N$-electron state is completely determined by the occupation
numbers of the First Order Reduced Density Operator (FORDO) and a set of one
particle \emph{Natural }General Spin Orbitals (NGSOs). The occupation numbers
in this case are equal to one or zero. In the AGP\ case the $2N$-electron
state is completely determined by a set of one particle \emph{Canonical
}General Spin Orbitals (CGSOs) and occupation numbers of the FORDO, but now
the occupation numbers can be any value in the range $\left[  0,1\right]  ,$
though they must be at least doubly degenerate. This theory \emph{explicitly
contains} correlation and generalizes the one-particle theory based on
uncorrelated Independent Particle States (IPS), which it contains as a special
case. The AGP\ model produces a simple pictorial interpretation in terms of
occupation numbers and orbitals, which totally \emph{determine} the
$2N$-electron state, and thus can be used as an effective and rigorous
chemical modelling tool.

In contrast to IPS's when using AGP's it is important to note the distinction between

\begin{itemize}
\item AGP-NGSO's, which together with a set of doubly degenerate occupation
numbers, determine an AGP-FORDO, that correspond to a \emph{set }of
AGP\ states and

\item AGP-CGSO's, which together with a set of doubly degenerate occupation
numbers, determine a \emph{unique} AGP state.
\end{itemize}

The CGSO's are the NGSO's of the FORDO of a given AGP state, but the NGSO's of
the FORDO\ of this AGP state are \emph{not }in general the CSGO's of this AGP\ state.

In this article we describe how the manifold of $2N$-electron AGP states can
be partitioned into equivalence classes, $\left[  D^{1}\right]  _{2N}$,
indexed by FORDO's $D^{1}$ that are at least doubly degenerate, and how the
elements of each class can be parameterized by a vector of angles
\begin{equation}
\mathbf{\theta\,}=\left(  0,\theta_{2},\cdots,\theta_{s}\right)
,\quad0<\theta_{j}\leq2\pi;\quad2\leq j\leq s,
\end{equation}
and a representative FORDO $D_{0}^{1}.$ As the occupation numbers $\left\{
N_{j}\quad1\leq j\leq s\right\}  $ of the $2N$-AGP\ state $\left\vert
g^{N}\right\rangle $ are in $1-1$ correspondence \cite{Colemann} with the
occupation numbers $\left\{  n_{j}\quad1\leq j\leq s\right\}  $ of the
$2$-electron geminal state $\left\vert g\right\rangle $ each representative
$D_{0}^{1}$ can, in turn, be parameterized by a vector of geminal occupation
numbers%
\begin{equation}
\mathbf{n}=\left(  n_{1},\cdots,n_{s}\right)  ,\quad0\leq n_{j}\leq
1;\quad1\leq j\leq s;\quad%
%TCIMACRO{\tsum \limits_{1\leq j\leq s}}%
%BeginExpansion
{\textstyle\sum\limits_{1\leq j\leq s}}
%EndExpansion
n_{j}=1
\end{equation}
and a representative set of NGSO's $\left\vert \mathbf{\varphi}\right\rangle
=\left\{  \left.  \left\vert \varphi_{j}\right\rangle \right\vert 1\leq
j\leq2s\right\}  $. The energy functional $\mathcal{E}\left(  g^{N}\right)  $
over the manifold of AGP\ states can thus be formulated as a functional
$\mathcal{E}\left(  \mathbf{n},\mathbf{\theta,}\left\vert \mathbf{\varphi
}\right\rangle \right)  $.

We show that the necessary conditions for the minimization of the energy wrt
$\left\vert \mathbf{\varphi}\right\rangle $, subject to the \emph{normality}
constraints of the GSO's $\left\{  \left.  \left\vert \varphi_{j}\right\rangle
\right\vert 1\leq j\leq2s\right\}  $, for fixed $\left(  \mathbf{n}%
,\mathbf{\theta}\right)  $ can be formulated in terms of a generalized Fock
operator $F\left(  \mathbf{n},\mathbf{\theta}\right)  $ (we display that is
sufficient to hold the normality constraint in order to obtain
orthonormality). By expressing the total energy in terms of the eigenvalues of
$F\left(  \mathbf{n},\mathbf{\theta}\right)  $ we deduce the same aufbau
principle as in the IPS\ case, i.e. in general the GSO's that are eigenvectors
of $F\left(  \mathbf{n},\mathbf{\theta}\right)  $ with with the lowest
eigenvalues should be associated with the largest occupation numbers. If the
occupation numbers are more than doubly degenerate or if the eigenvalues of
$F\left(  \mathbf{n},\mathbf{\theta}\right)  $ are degenerate than multiple
assignments may have to be examined. After all the variables have been
optimized one can utilize the Fock operator $F\left(  \mathbf{n}%
,\mathbf{\theta}\right)  $ as a first order approximation to the Hamiltonian
in the one-electron propagator and again associate the eigenvalues $\left\{
\left.  \varepsilon_{j}\right\vert 1\leq j\leq2s\right\}  $ with ionization potentials.

Unlike the IPS\ case the phases of the eigenvectors of $F\left(
\mathbf{n},\mathbf{\theta}\right)  $ \emph{are }important as they effect the
total energy and an optimal value of $\mathbf{\theta}$ must be found. These
phases are not determined by the diagonalization of $F\left(  \mathbf{n}%
,\mathbf{\theta}\right)  $ and have to be evaluated by an independent
variational equation. This leads to a Density Matrix Functional (DMF),
$\mathcal{F}$, defined on the set of equivalence classes $\left\{  \left[
D^{1}\right]  \right\}  $ by\
\begin{equation}
\mathcal{F}\left(  D^{1}\right)  =\min_{\mathbf{\theta}}\left\{
\mathcal{E}\left(  \mathbf{n},\mathbf{\theta,}\left\vert \mathbf{\varphi
}\right\rangle \right)  \right\}  ,
\end{equation}
which can also be utilized in Density Matrix Functional Theory (DMFT) as an
explicitly known example of a DMF.

The geminal occupation numbers $\left\{  n_{j}\quad1\leq j\leq s\right\}  $
are optimized using an iterative matrix procedure that was introduced
\cite{weiner,ortiz,ohrn,jensen} many years ago, which we show automatically
enforces the normalization and positivity conditions.

The variational procedure can be schematically displayed as
\begin{align}
&  \left(  \mathbf{n}_{i},\mathbf{\theta}_{i}\mathbf{,}\left\vert
\mathbf{\varphi}_{i}\right\rangle \right)  \overset{\text{ }}{\qquad
\qquad\qquad\underrightarrow{\text{iterate on the Fock Operator}}\qquad\left(
\mathbf{n}_{i},\mathbf{\theta}_{i}\mathbf{,}\left\vert \mathbf{\varphi}%
_{i+1}\left(  \mathbf{n}_{i},\mathbf{\theta}_{i}\right)  \right\rangle
\right)  }\nonumber\\
&  \qquad\left\uparrow
\begin{array}
[c]{c}%
\\
i\rightarrow i+1\\
\end{array}
\right.  \qquad\qquad\qquad\qquad\qquad\qquad\qquad\qquad\qquad
\left\downarrow
\begin{array}
[c]{c}%
\\
\text{minimize wrt }\mathbf{\theta}\\
\end{array}
\right. \nonumber\\
&  \left(  \mathbf{n}_{i+1},\mathbf{\theta}_{i+1}\mathbf{,}\left\vert
\mathbf{\varphi}_{i+1}\left(  \mathbf{n}_{i},\mathbf{\theta}_{i+1}\right)
\right\rangle \right)  \underleftarrow{\qquad\text{minimize wrt }%
\mathbf{n}\qquad\qquad}\left(  \mathbf{n}_{i},\mathbf{\theta}_{i+1}%
\mathbf{,}\left\vert \mathbf{\varphi}_{i+1}\left(  \mathbf{n}_{i}%
,\mathbf{\theta}_{i+1}\right)  \right\rangle \right)  .
\end{align}
The GSO's and $\mathbf{\theta}$ that are self consistent solutions to the
above scheme determine the CGSO's and together with the self consistent
occupation numbers $\mathbf{n}$ determine the optimal geminal and
$2N$-electron AGP state.

The one electron model that we obtain describes molecules by weighted PD's,
each determined by an occupation number and a phase optimized GSO associated
with an one particle energy, in an analogous fashion to the IPS case, but also
contains correlation effects. Pictorial manipulations of these PD's (taking
into account relaxation of the one particle energies and optimal phases)
allows one to analyze chemical processes and structures in a rigorously
meaningful fashion, as one is assured that the picture produced corresponds to
multi electron states.

The structure of this article is as follows: in section \ref{AGPstate} we
describe the canonical form of a geminal, and the AGP\ state, in terms of
geminal occupation numbers $\left\{  \left.  n_{j}\right\vert 1\leq j\leq
s\right\}  $ and CGSOs $\left\{  \left.  \left\vert \psi_{j}\right\rangle
\right\vert 1\leq j\leq2s\right\}  .$ Together with the AGP-FORDO occupation
numbers $\left\{  \left.  N_{j}\right\vert 1\leq j\leq s\right\}  ,$ which are
in $1-1$ correspondence with the geminal occupation numbers, these CGSOs
determine both the AGP\ state and the AGP-FORDO. In section \ref{transform} we
examine the $1-1$ transformations between the AGP $\left\{  \left.
N_{j}\right\vert 1\leq j\leq s\right\}  $ and geminal $\left\{  \left.
n_{j}\right\vert 1\leq j\leq s\right\}  $ occupation numbers. In section
\ref{Energy} we discuss the energy expressions in terms of the occupation
numbers $\left\{  \left.  n_{j}\right\vert 1\leq j\leq s\right\}  $, CGSOs
$\left\{  \left.  \left\vert \psi_{j}\right\rangle \right\vert 1\leq
j\leq2s\right\}  $ and phases $\left\{  \left.  \theta_{j}\right\vert 1\leq
j\leq s\right\}  .$ In section \ref{Evariation} we describe the necessary
conditions needed to minimize the energy over a manifold of AGP\ states. In
section \ref{Fock} we derive a generalized Fock operator that produces the
derivative of the energy functional with respect to NGSO's for fixed FORDO
occupation numbers and phases. In section \ref{aufbau} We analyze the total
energy in terms of the eigenvalues of this Fock operator and show that an
aufbau principle similar to the IPS\ case is effective in constructing
multiparticle states from the eigenvectors of the Fock operator. In section
\ref{theta} we solve the equations that produces the optimal phases
$\mathbf{\theta.}$ In section \ref{gemocc} we describe the iterative procedure
that determines the optimal geminal occupation numbers $\mathbf{n}$. In
\ref{procedure} we form a combined procedure that produces an energy optimized
AGP state and we conclude in section \ref{Discussion} with a review of the
material presented and a discussion of its implications.

\section{The $2N$-electron AGP state\label{AGPstate}}

In this section we review the canonical form of a geminal and AGP state and
their associated FORDO's. We also characterize equivalence classes of AGP
states that have the same FORDO.

\subsection{Canonical Form}

The $2N$-electron AGP state, $\left\vert g^{N}\right\rangle $, is the $N^{th}$
antisymmetric power of the $2$-electron state described by the geminal
\begin{equation}
\left\vert g\right\rangle =\sum_{1\leq j\leq s}n_{j}^{\frac{1}{2}}\left\vert
\psi_{j}\psi_{j+s}\right\rangle , \label{geminal}%
\end{equation}
where $\left\{  \left.  n_{j}^{\frac{1}{2}}\right\vert 1\leq j\leq s\right\}
$ are the positive square roots of $\left\{  \left.  n_{j}\right\vert 1\leq
j\leq s\right\}  ,$ and is given by \cite{weiner&ortiz2004}
\begin{equation}
\left\vert g^{N}\right\rangle =\left(  S_{N}\right)  ^{-\frac{1}{2}}%
\sum_{1\leq j_{1}<\cdots<j_{N}\leq s}n_{j_{1}}^{\frac{1}{2}}\cdots n_{j_{s}%
}^{\frac{1}{2}}\left\vert \psi_{j_{1}}\psi_{j_{1}+s}\cdots\psi_{j_{N}}%
\psi_{j_{N}+s}\right\rangle . \label{agp}%
\end{equation}
The normalization factor, $\left(  S_{N}\right)  ^{-\frac{1}{2}}$, is given in
terms of the elementary symmetric function
\begin{equation}
S_{N}=\sum_{1\leq j_{1}<\cdots<j_{N}\leq s}n_{j_{1}}\cdots n_{j_{N}},
\end{equation}
and the geminals $\left\vert g\right\rangle $ represent normalized two
electron states so that
\begin{equation}
\sum_{1\leq j\leq s}n_{j}=1.
\end{equation}
It is evident from the expansion eq. $\left(  \ref{geminal}\right)  $ that the
geminal and geminal power state are at the least invariant to the group
\begin{equation}
SU\left(  2\right)  ^{s}=\overset{s\text{-factors}}{\overbrace{\,SU\left(
2\right)  \times\cdots\times SU\left(  2\right)  }},
\end{equation}
that is formed by $2\times2$ \emph{special} unitary transformations of the
paired CGSO's $\left\{  \left.  \left\vert \psi_{j}\right\rangle ,\left\vert
\psi_{j+s}\right\rangle \right\vert 1\leq j\leq s\right\}  ,$ while if
degeneracies exist amongst the $\left\{  \left.  n_{j}\right\vert 1\leq j\leq
s\right\}  $ this invariance group is larger \cite{Weiner2004}. Hence the
canonical CSGO's are unique up to transformations belonging this invariance group.

The CGSO's $\left\{  \left.  \left\vert \psi_{j}\right\rangle \right\vert
1\leq j\leq2s\right\}  $ can be expanded in terms of a spin orbitals basis
$\left\{  \left.  \left\vert \varphi_{j\alpha}\right\rangle ,\left\vert
\varphi_{j\beta}\right\rangle \right\vert 1\leq j\leq s\right\}  $ of one
electron state space $\mathcal{H}^{1}$ as
\begin{equation}
\left\vert \psi_{j}\right\rangle =\sum_{1\leq j\leq s}\left\{  a_{j}\left\vert
\varphi_{j\alpha}\right\rangle +b_{j}\left\vert \varphi_{j\beta}\right\rangle
\right\}  ,
\end{equation}
leading to one electron GSO wavefunctions defined on the spectrum of position
and spin operators $\left(  \text{see Appendix \ref{basis}}\right)  $ as
\begin{equation}
\psi_{j}\left(  \mathbf{r,}\zeta\right)  =\left\langle \left.  \mathbf{r,}%
\zeta\right\vert \psi_{j}\right\rangle .
\end{equation}
It should be noted that if expansion eq. $\left(  \ref{geminal}\right)  $ is
used to define the canonical expansion it is usually \emph{not possible} to
find a set of CSGO's such that
\begin{equation}
\psi_{j+s}\left(  \mathbf{r,}\zeta\right)  =\psi_{j}\left(  \mathbf{r,}%
\zeta\right)  ^{\ast}.
\end{equation}


\subsection{First Order Reduced Density Operators}

The states $\left\vert g\right\rangle $ and $\left\vert g^{N}\right\rangle $
define FORDO's given by \cite{Coleman}
\begin{equation}
D^{1}\left(  g\right)  =\sum_{1\leq j\leq s}n_{j}\left\{  \left\vert \psi
_{j}\right\rangle \left\langle \psi_{j}\right\vert +\left\vert \psi
_{j+s}\right\rangle \left\langle \psi_{j+s}\right\vert \right\}  \label{d1g}%
\end{equation}
and
\begin{equation}
D^{1}\left(  g^{N}\right)  =\sum_{1\leq j\leq s}N_{j}\left\{  \left\vert
\psi_{j}\right\rangle \left\langle \psi_{j}\right\vert +\left\vert \psi
_{j+s}\right\rangle \left\langle \psi_{j+s}\right\vert \right\}  \label{dng}%
\end{equation}
respectively. It is clear from the expansions eqs. $\left(  \ref{d1g}\right)
$ and $\left(  \ref{dng}\right)  $ that the CGSO's $\left\{  \left.
\left\vert \psi_{j}\right\rangle \right\vert 1\leq j\leq r\right\}  $ are
simultaneously the Natural General Spin Orbitals (NGSO's) of $D^{1}\left(
g\right)  $ and $D^{1}\left(  g^{N}\right)  $. However it is very important to
note that the NGSO's of $D^{1}\left(  g\right)  $ and $D^{1}\left(
g^{N}\right)  $ \emph{are not }in general the CGSO's of $\left\vert
g\right\rangle $. This is a ramification of

\begin{enumerate}
\item the general invariance of FORDO's to the multiplication of their NGSO's
by complex phase factors
\begin{equation}
D^{1}=\sum_{1\leq j\leq r}\nu_{j}\left\{  \left\vert \phi_{j}\right\rangle
\left\langle \phi_{j}\right\vert \right\}  =\sum_{1\leq j\leq r}\nu
_{j}\left\{  e^{i\theta_{j}}\left\vert \phi_{j}\right\rangle \left\langle
\phi_{j}\right\vert e^{-i\theta_{j}}\right\}  ,
\end{equation}
where $\left\{  \left.  \nu_{j}\right\vert 1\leq j\leq s\right\}  $ are
occupation numbers and

\item the observation that geminals \emph{are not }invariant to such phase
changes of their CGSO's.
\end{enumerate}

One can form equivalence classes $\left[  D^{1}\left(  g\left(  \mathbf{0}%
\right)  \right)  \right]  $ of geminals $\left\{  \left\vert g\left(
\mathbf{\theta}\right)  \right\rangle \right\}  $ that have the same FORDO
$D^{1}\left(  g\left(  \mathbf{0}\right)  \right)  $ and equivalence classes
$\left[  D^{1}\left(  g^{N}\left(  \mathbf{0}\right)  \right)  \right]  $ of
$2N$-electron AGP states $\left\{  \left\vert g^{N}\left(  \mathbf{\theta
}\right)  \right\rangle \right\}  $, that have the same FORDO $D^{1}\left(
g^{N}\left(  \mathbf{0}\right)  \right)  $. For a fixed $N$ these classes
correspond to each other in a bijective manner \cite{coleman} and can be
parameterized by the angles $\mathbf{\theta}$, where
\begin{align}
\left\vert g\left(  \mathbf{\theta}\right)  \right\rangle  &  =\sum_{1\leq
j\leq s}n_{j}^{\frac{1}{2}}e^{i\theta_{j}}\left\vert \psi_{j}\psi
_{j+s}\right\rangle =\sum_{1\leq j\leq s}n_{j}^{\frac{1}{2}}\left\vert
e^{i\varphi_{j}}\psi_{j}e^{i\varphi_{j+s}}\psi_{j+s}\right\rangle ;\nonumber\\
&  \qquad\qquad\qquad\qquad\qquad\qquad\qquad\theta_{1}=0,0<\theta_{j}\leq
2\pi,2\leq j\leq s,
\end{align}
where $\mathbf{\theta=}\left(  0,\theta_{2},\cdots,\theta_{s}\right)  $ and
$\varphi_{j}+\varphi_{j+s}=\theta_{j}$. Choosing $\theta_{1}=0$ eliminates the
overall phase factor indeterminacy of $\left\vert g\left(  \mathbf{\theta
}\right)  \right\rangle $. The individual phase factors $\left\{  \left.
\varphi_{j}\right\vert 1\leq j\leq s\right\}  $ are not unique, though their
sums $\left\{  \left.  \theta_{j}\right\vert 1\leq j\leq s\right\}  $ are,
allowing one to obtain a unique representative set $\left\{  \left.
\varphi_{j}\right\vert 1\leq j\leq s\right\}  $ by defining
\begin{equation}
\varphi_{j}=\frac{\theta_{j}}{2},\varphi_{j+s}=\frac{\theta_{j}}{2};\quad1\leq
j\leq s.
\end{equation}
Thus one has the correspondence
\begin{equation}
\left\{  \left\vert g\left(  \mathbf{\theta}\right)  \right\rangle
\longleftrightarrow\left\vert g^{N}\left(  \mathbf{\theta}\right)
\right\rangle ;\quad\theta_{1}=0;\,0<\theta_{j}\leq2\pi,2\leq j\leq s\right\}
\longleftrightarrow\left[  D^{1}\left(  g^{N}\left(  \mathbf{0}\right)
\right)  \right]  .
\end{equation}
The elements of the class $\left[  D^{2}\left(  g^{N}\left(  \mathbf{0}%
\right)  \right)  \right]  $ can also be expressed in canonical form as
\begin{equation}
\sum_{1\leq j\leq s}n_{j}^{\frac{1}{2}}\left\vert \chi_{j}\chi_{j+s}%
\right\rangle ,
\end{equation}
where $\left\vert \chi_{j}\right\rangle =\left\vert e^{i\varphi_{j}}\psi
_{j}\right\rangle $. The AGP\ states $\left\vert g^{N}\left(  \mathbf{\theta
}\right)  \right\rangle $ are clearly distinct as they have different Second
Order Reduced Density Operators (SORDO's), $D^{2}\left(  g^{N}\left(
\mathbf{\theta}\right)  \right)  $, given by \cite{Coleman}
\begin{align}
&  D^{2}\left(  g^{N}\left(  \mathbf{\theta}\right)  \right)  =\frac{1}{S_{N}%
}\left\{  \sum_{1\leq j\leq s}n_{j}\left\vert \psi_{j}\psi_{j+s}\right\rangle
\left\langle \psi_{j}\psi_{j+s}\right\vert \right.  +\nonumber\\
&  \left.  \sum_{1\leq j,k\leq s}n_{j}^{\frac{1}{2}}n_{k}^{\frac{1}{2}%
}e^{i\left(  \theta_{j}-\theta_{k}\right)  }S_{N-1}\left(  \jmath k\right)
\left\vert \psi_{j}\psi_{j+s}\right\rangle \left\langle \psi_{k}\psi
_{k+s}\right\vert +\sum_{1\leq j<k\leq s}n_{j}n_{k}S_{N-2}\left(  \jmath
k\right)  \left\vert \psi_{j}\psi_{k}\right\rangle \left\langle \psi_{j}%
\psi_{k}\right\vert \right\}  , \label{D2}%
\end{align}
where
\begin{equation}
S_{N-1}\left(  \jmath k\right)  =\sum_{\substack{1\leq j_{1}<\cdots
<j_{N-1}\leq s\\j,k\notin\left\{  j_{1},\cdots,j_{N-1}\right\}  }}n_{j_{1}%
}\cdots n_{j_{N-1}}%
\end{equation}
and
\begin{equation}
S_{N-2}\left(  \jmath k\right)  =\sum_{\substack{1\leq j_{1}<\cdots
<j_{N-2}\leq s\\j,k\notin\left\{  j_{1},\cdots,j_{N-2}\right\}  }}n_{j_{1}%
}\cdots n_{j_{N-2}}.
\end{equation}
Note that $S_{N-1}\left(  \jmath\jmath\right)  =S_{N-2}\left(  \jmath
\jmath\right)  =0.$ The SORDO eq. $\left(  \ref{D2}\right)  $ can be expanded
in terms of $\mathbf{\theta}$ dependent CGSO's $\left\{  \left.  \left\vert
e^{i\varphi_{j}}\chi_{j}\right\rangle ,\left\vert e^{i\varphi_{j+s}}\chi
_{j+s}\right\rangle \right\vert 1\leq j\leq s\right\}  ,$ where $\left\{
\left.  \varphi_{j}+\varphi_{j+s}=\theta_{j}\right\vert 1\leq j\leq s\right\}
$ as
\begin{align}
&  D^{2}\left(  g^{N}\left(  \mathbf{\theta}\right)  \right)  =\frac{1}{S_{N}%
}\left\{  \sum_{1\leq j\leq s}n_{j}S_{N-1}\left(  \jmath\right)  \left\vert
e^{i\varphi_{j}}\chi_{j}e^{i\varphi_{j+s}}\chi_{j+s}\right\rangle \left\langle
e^{i\varphi_{j}}\chi_{j}e^{i\varphi_{j+s}}\chi_{j+s}\right\vert \right.
\nonumber\\
&  \qquad\qquad\qquad\qquad+\sum_{1\leq j,k\leq s}n_{j}^{\frac{1}{2}}%
n_{k}^{\frac{1}{2}}S_{N-1}\left(  \jmath k\right)  \left\vert e^{i\varphi_{j}%
}\chi_{j}e^{i\varphi_{j+s}}\chi_{j+s}\right\rangle \left\langle e^{i\varphi
_{k}}\chi_{k}e^{i\varphi_{k+s}}\chi_{k+s}\right\vert \nonumber\\
&  \qquad\qquad\qquad\qquad\left.  +\sum_{1\leq j<k\leq s}n_{j}n_{k}%
S_{N-2}\left(  \jmath k\right)  \left\vert e^{i\varphi_{j}}\chi_{j}%
e^{i\varphi_{k+s}}\chi_{k+s}\right\rangle \left\langle e^{i\varphi_{j}}%
\chi_{j}e^{i\varphi_{k+s}}\chi_{k+s}\right\vert \right\}  , \label{en2}%
\end{align}
giving the explicitly phase dependent expression
\begin{align}
&  D^{2}\left(  g^{N}\left(  \mathbf{\theta}\right)  \right)  =\frac{1}{S_{N}%
}\left\{  \sum_{1\leq j\leq s}n_{j}S_{N-1}\left(  \jmath\right)  \left\vert
\chi_{j}\chi_{j+s}\right\rangle \left\langle \chi_{j}\chi_{j+s}\right\vert
+\sum_{1\leq j,k\leq s}n_{j}^{\frac{1}{2}}n_{k}^{\frac{1}{2}}S_{N-1}\left(
\jmath k\right)  e^{i\left(  \theta_{j}-\theta_{k}\right)  }\left\vert
\chi_{j}\chi_{j+s}\right\rangle \left\langle \chi_{k}\chi_{k+s}\right\vert
\right. \nonumber\\
&  \left.  +\sum_{1\leq j<k\leq s}n_{j}n_{k}S_{N-2}\left(  \jmath k\right)
\left\{  \left\vert \chi_{j}\chi_{k}\right\rangle \left\langle \chi_{j}%
\chi_{k}\right\vert +\left\vert \chi_{j+s}\chi_{k}\right\rangle \left\langle
\chi_{j}\chi_{k+s}\right\vert +\left\vert \chi_{j}\chi_{k+s}\right\rangle
\left\langle \chi_{j}\chi_{k+s}\right\vert +\left\vert \chi_{j+s}\chi
_{k+s}\right\rangle \left\langle \chi_{j+s}\chi_{k+s}\right\vert \right\}
\right\}  .
\end{align}


\section{The correspondence between AGP and Geminal Occupation
Numbers\label{transform}}

Any set of $s$ positive numbers $\left\{  \left.  N_{j}\right\vert 1\leq j\leq
s\right\}  $ that have values between $0$ and $1$ and add up to $N,$ can be
interpreted as the occupation numbers of a FORDO that is at least doubly
degenerate and can be associated with an AGP\ state. These numbers are in
$1-1$ correspondence with another set of positive numbers $\left\{  \left.
n_{j}\right\vert 1\leq j\leq s\right\}  ,$ that add up to $1$, which can be
interpreted as the occupation numbers of a FORDO\ of a geminal. The geminal
occupation numbers $\left\{  \left.  n_{j}\right\vert 1\leq j\leq s\right\}  $
together with an arbitrary set of $2s$ orthonormal one electron Canonical
General Spin Orbitals (CGSO's) then \emph{uniquely }determine a $2N$ electron
AGP state. Any state that has a FORDO that has at least doubly degenerate
occupation numbers is covered by a set of AGP\ states i.e. there exists
AGP\ states that have the same FORDO. In most applications e.g. expectation
values of observables, it is sufficient to know the $2N$ electron state
occupation numbers in terms of the geminal occupation numbers, which serve as
a parameterization of the state occupation numbers, and we review this
relationship first. The inverse function i.e. geminal occupation numbers in
terms of state occupation numbers are significant in designing possible
AGP\ states and we present a numerical construction of this function.

\subsection{Transformation of Geminal to AGP Occupation Numbers}

The vector of occupation numbers $\mathbf{N}=\left(  N_{1},\cdots
,N_{s}\right)  $ of an AGP\ state $\left\vert g^{N}\right\rangle $ can be
expressed as a bijective function, $\mathcal{N}\left(  \mathbf{n}\right)  $,
of the geminal occupation numbers $\mathbf{n}=\left(  n_{1},\cdots
,n_{s}\right)  ,$ where
\begin{equation}
\mathcal{N}:\mathfrak{D}_{N}\subset\mathbb{R}^{s}\rightarrow\mathfrak{D}%
_{n}\subset\mathbb{R}^{s}.
\end{equation}
The domain $\mathfrak{D}_{N}$ and range $\mathfrak{D}_{n}$ are given by%
\begin{align}
\mathfrak{D}_{N}  &  =\left(  N_{1},\cdots,N_{s}\right)  ;\quad0\leq N_{j}%
\leq1,\quad1\leq j\leq s\nonumber\\
\sum_{1\leq j\leq s}N_{j}  &  =N
\end{align}
and
\begin{align}
\mathfrak{D}_{n}  &  =\left(  n_{1},\cdots,n_{s}\right)  ;\quad n_{j}%
\geq0,\quad1\leq j\leq s\nonumber\\
\sum_{1\leq j\leq s}n_{j}  &  =1,
\end{align}
where
\begin{equation}
N_{j}=\mathcal{N}_{j}\left(  \mathbf{n}\right)  =n_{j}\frac{S_{N-1}\left(
\jmath\right)  }{S_{N}};\quad1\leq j\leq s
\end{equation}
and
\begin{equation}
S_{N-1}\left(  \jmath\right)  =\sum_{\substack{1\leq j_{1}<\cdots<j_{N-1}\leq
s\\j\notin\left\{  j_{1},\cdots,j_{N-1}\right\}  }}n_{j_{1}}\cdots n_{j_{N-1}%
}.
\end{equation}


\subsection{Transformation of AGP to Geminal Occupation Numbers}

The inverse function $\mathcal{N}^{-1}\left(  \mathbf{N}\right)  $ is not
explicitly known, but one can obtain a linear approximation that suffices in
practical applications to numerical determine values of $\mathbf{n}$ from
known values of $\mathbf{N}$, i.e. the determination of $\mathbf{n}_{0}%
+\delta\mathbf{n}$ from $\mathbf{N}_{0}+\delta\mathbf{N}$, when $\mathbf{N}%
_{0}$ is a known initial vector of occupation numbers of the $2N$-electron AGP
state, e.g. an IPS,\ where
\begin{equation}
\mathbf{N}_{0}=\left(  \overset{N}{\overbrace{1,\cdots,1}},\overset
{s-N}{\overbrace{0,\cdots,0}}\right)
\end{equation}
and
\begin{equation}
\mathbf{n}_{0}=\left(  \overset{N}{\overbrace{\frac{1}{N},\cdots,\frac{1}{N}}%
},\overset{s-N}{\overbrace{0,\cdots,0}}\right)  .
\end{equation}
The Taylor series expansion about $\mathbf{N}_{0}$ of $\mathcal{N}^{-1}$ is
given by%
\begin{equation}
\mathcal{N}^{-1}\left(  \mathbf{N}_{0}+\delta\mathbf{N}\right)  =\mathcal{N}%
^{-1}\left(  \mathbf{N}_{0}\right)  +d^{1}\mathcal{N}^{-1}\left(
\mathbf{N}_{0}\right)  \left(  \delta\mathbf{n}\right)  +\cdots,
\end{equation}
where
\begin{equation}
d^{1}\mathcal{N}^{-1}\left(  \mathbf{N}_{0}\right)  =\left[  d^{1}%
\mathcal{N}\left(  \mathbf{n}_{0}\right)  \right]  ^{-1}%
\end{equation}
and the Jacobian, $d^{1}\mathcal{N}\left(  \mathbf{n}_{0}\right)  $, of the
transformation $\mathcal{N}$ at the point $\mathbf{n}_{0}$ is given by the
matrix elements
\begin{align}
\left[  d^{1}\mathcal{N}\left(  \mathbf{n}_{0}\right)  \right]  _{jk}  &
=\frac{\partial\mathcal{N}_{j}\left(  \mathbf{n}_{0}\right)  }{\partial n_{k}%
}=\frac{\partial}{\partial n_{k}}\left\{  n_{j}\frac{S_{N-1}\left(
\jmath\right)  }{S_{N}}\right\} \nonumber\\
&  =\frac{1}{S_{N}}\left(  \delta_{jk}S_{N-1}\left(  \jmath\right)
+n_{j}\frac{\partial S_{N-1}\left(  \jmath\right)  }{\partial n_{k}}%
-n_{j}\frac{S_{N-1}\left(  \jmath\right)  }{S_{N}}\frac{\partial S_{N}%
}{\partial n_{k}}\right) \nonumber\\
&  =\frac{1}{S_{N}}\left(  \delta_{jk}S_{N-1}\left(  \jmath\right)
+n_{j}S_{N-2}\left(  \jmath k\right)  -n_{j}\frac{1}{S_{N}}S_{N-1}\left(
\jmath\right)  S_{N-1}\left(  k\right)  \right)  .
\end{align}
The Jacobian of the inverse transformation $\mathcal{N}^{-1}\left(
\mathbf{N}_{0}\right)  $ at the point $\mathbf{N}_{0}$ is given by the inverse
of the matrix $\left\{  \left.  \left[  d^{1}\mathcal{N}\left(  \mathbf{n}%
_{0}\right)  \right]  _{jk}\right\vert 1\leq j,k\leq s\right\}  $ i.e.
\begin{equation}
d^{1}\mathcal{N}^{-1}\left(  \mathbf{N}_{0}\right)  =\left[  d^{1}%
\mathcal{N}\left(  \mathbf{n}_{0}\right)  \right]  ^{-1}%
\end{equation}
and the incremental change, $\delta\mathbf{n}$, of $\mathbf{n}_{0}$ is
\begin{equation}
d^{1}\mathcal{N}^{-1}\left(  \mathbf{N}_{0}\right)  \left(  \delta
\mathbf{N}\right)  _{j}=\sum_{1\leq k\leq s}\left[  d^{1}\mathcal{N}\left(
\mathbf{n}_{0}\right)  \right]  _{jk}^{-1}\left(  \delta\mathbf{N}\right)
_{k}.
\end{equation}
We can then construct a piecewise linear fit to $\mathcal{N}^{-1}$ as%
\begin{equation}
\mathbf{n}\simeq\mathbf{\tilde{n}}_{m}=\mathcal{N}^{-1}\left(  \mathbf{N}%
\right)  =\mathbf{n}_{0}+\sum_{1\leq k\leq m}\left[  d^{1}\mathcal{N}\left(
\mathbf{n}_{0}+\sum_{1\leq j\leq k-1}\left[  d^{1}\mathcal{N}\left(
\mathbf{n}_{0}\right)  \right]  ^{-1}\left(  \delta\mathbf{N}_{j}\right)
\right)  \right]  ^{-1}\left(  \delta\mathbf{N}_{k}\right)  ,
\end{equation}
where
\begin{equation}
\mathbf{N=N}_{0}+\sum_{1\leq j\leq k-1}\delta\mathbf{N}_{j}%
\end{equation}
and
\begin{equation}
\sum_{1\leq k\leq s}\delta\mathbf{N}_{jk}=0.
\end{equation}
See Appendix \label{linfit} for an example. The accuracy of the fit can be
checked by evaluating
\begin{equation}
\mathbf{\Delta}_{m}=\mathcal{N}\left(  \mathbf{\tilde{n}}_{m}\right)
-\mathbf{N}_{m}%
\end{equation}
and noting if
\begin{equation}
\left\vert \mathbf{\Delta}_{mj}\right\vert \leq\varkappa_{j};\quad1\leq j\leq
s \label{con}%
\end{equation}
If the error $\left(  \left\vert \mathbf{\Delta}_{m1}\right\vert
,\cdots,\left\vert \mathbf{\Delta}_{ms}\right\vert \right)  $, i.e. is too big
i.e. eq. $\left(  \ref{con}\right)  $ is not satisfied one can then decrease
the size of $\delta\mathbf{N}_{m},$ etc.

\section{Energy Functionals \label{Energy}}

In this section we describe the energy of a $2N$-electron AGP state in three
different ways in terms of

\begin{description}
\item
\begin{enumerate}
\item CGSO's and geminal occupation numbers: 

\item A density matrix functional $\mathcal{F}\left(  D^{1}\right)  $ defined
on the equivalence classes $\left[  D^{1}\left(  g^{N}\left(  \mathbf{0}%
\right)  \right)  \right]  .$ The functional $\mathcal{F}\left(  D^{1}\right)
$ is significant in two aspects, it can serve as a rigorous explicitly known
model functional in Density Matrix Functional Theory (DMFT) \cite{DMFT} and it
clealry demostrtaes the difference between CGSO's and NGSO's.

\item Group parameters introduced in \cite{weiner2004} that describe the
manifold of agp\ states.
\end{enumerate}
\end{description}

\subsection{Energy in terms of Geminal Occupation Numbers and CGSO's}

The energy, $\mathcal{E}\left(  \mathbf{n,\psi}\right)  ,$ of a $2N$-electron
AGP\ state for a molecule at a fixed nuclear configuration $\left(
\mathbf{R}_{1},\cdots,\mathbf{R}_{N_{nuc}}\right)  $ in terms of the CGSO's
amplitudes $\left\{  \left.  \left\langle \left.  \mathbf{r,\xi}\right\vert
\psi_{j}\right\rangle \right\vert 1\leq j\leq r\right\}  $ and geminal
occupation numbers $\mathbf{n}$ is
\begin{align}
\mathcal{E}\left(  \mathbf{n,\psi}\right)   &  =\frac{\mathcal{H}\left(
\mathbf{n,\psi}\right)  }{S_{N}\left(  \mathbf{n}\right)  }\nonumber\\
&  =\frac{1}{S_{N}\left(  \mathbf{n}\right)  }\left\{  \sum_{1\leq j\leq
s}n_{j}S_{N-1}\left(  \jmath\right)  \left\{  \left\langle \psi_{j}\left\vert
h_{1}\psi_{j}\right.  \right\rangle +\left\langle \psi_{j+s}\left\vert
h_{1}\psi_{j+s}\right.  \right\rangle +\left\langle \psi_{j}\psi_{j+s}%
||\psi_{j}\psi_{j+s}\right\rangle \right\}  \right. \nonumber\\
&  +\sum_{1\leq j,k\leq s}n_{j}^{\frac{1}{2}}n_{k}^{\frac{1}{2}}S_{N-1}\left(
\jmath k\right)  \left\langle \psi_{j}\psi_{j+s}||\psi_{k}\psi_{k+s}%
\right\rangle \left.  +\sum_{1\leq j<k\leq s}n_{i}n_{j}S_{N-2}\left(  \jmath
k\right)  \left\langle \psi_{j}\psi_{k}|||\psi_{j}\psi_{k}\right\rangle
\right\}  . \label{energy}%
\end{align}
$\lceil$Note: Mathematically the energy functional should really be denoted as
$\mathcal{E}\left(  \mathbf{n,\psi}^{\ast}\mathbf{,\psi}\right)  $ as
$\mathbf{\psi}^{\ast}$ and $\mathbf{\psi}$ are treated as independent
variables even though $\mathbf{\psi}$ determines $\mathbf{\psi}^{\ast}%
$.$\rfloor$ The one-electron matrix elements $\left\{  \left\langle \psi
_{j}\left\vert h_{1}\psi_{k}\right.  \right\rangle \right\}  $ wrt the CGSO's
are
\begin{equation}
\left\langle \psi_{j}\left\vert h_{1}\psi_{k}\right.  \right\rangle =-\int
\psi_{j}^{\ast}\left(  \mathbf{r,\xi}\right)  \left\{  \sum_{_{1\leq l\leq
N_{nuc}}}\frac{Z_{l}}{\left\Vert \mathbf{R}_{l}-\mathbf{r}\right\Vert }%
+\nabla_{\mathbf{r}}^{2}\right\}  \psi_{k}\left(  \mathbf{r,\xi}\right)
d\mathbf{r}d\xi;\quad1\leq j,k\leq2s, \label{h1}%
\end{equation}
the antisymmetric two-electron matrix elements $\left\{  \left\langle
jk||lm\right\rangle \right\}  $
\begin{align}
\left\langle jk||lm\right\rangle  &  =\int\left\{  \psi_{j}^{\ast}\left(
\mathbf{r}_{1}\mathbf{,\xi}_{1}\right)  \psi_{k}^{\ast}\left(  \mathbf{r}%
_{2}\mathbf{,\xi}_{2}\right)  \frac{1}{\left\Vert \mathbf{r}_{1}%
-\mathbf{r}_{2}\right\Vert }\psi_{l}\left(  \mathbf{r}_{1}\mathbf{,\xi}%
_{1}\right)  \psi_{m}\left(  \mathbf{r}_{2}\mathbf{,\xi}_{2}\right)  \right.
\nonumber\\
&  -\left.  \psi_{j}^{\ast}\left(  \mathbf{r}_{1}\mathbf{,\xi}_{1}\right)
\psi_{k}^{\ast}\left(  \mathbf{r}_{2}\mathbf{,\xi}_{2}\right)  \frac
{1}{\left\Vert \mathbf{r}_{1}-\mathbf{r}_{2}\right\Vert }\psi_{l}\left(
\mathbf{r}_{2}\mathbf{,\xi}_{1}\right)  \psi_{m}\left(  \mathbf{r}%
_{1}\mathbf{,\xi}_{2}\right)  \right\}  d\mathbf{r}_{1}d\mathbf{r}_{2}d\xi
_{1}d\xi_{2}\nonumber\\
&  \left.  {}\right.  \qquad\qquad\qquad\qquad\qquad\qquad\qquad\qquad
\qquad1\leq j,k\leq2s;1\leq l,m\leq2s \label{h2}%
\end{align}
and $\left\{  \left\langle jk|||jk\right\rangle \right\}  $
\begin{align}
\left\langle \psi_{j}\psi_{k}|||\psi_{j}\psi_{k}\right\rangle  &
=\left\langle jk||jk\right\rangle +\left\langle j\,k+s||j\,k+s\right\rangle
+\left\langle j+s\,k||j+s\,k\right\rangle +\left\langle
j+s\,k+s||j+s\,k+s\right\rangle \nonumber\\
&  \left.  {}\right.  \qquad\qquad\qquad\qquad\qquad\qquad\qquad\qquad
\qquad\qquad\qquad\qquad\qquad\qquad1\leq j,k\leq s.
\end{align}


\subsection{Energy in terms of CGSO's Phases and FORDO's}

In terms of $\theta$ dependent CGSO's $\left\{  \left\vert e^{i\varphi_{j}%
}\psi_{j}\right\rangle ;1\leq j\leq2s\right\}  $ defined in eq. $\left(
\ref{en2}\right)  $ we obtain an energy expression factored through the
equivalence classes $\left[  D^{1}\left(  g^{N}\left(  \mathbf{0}\right)
\right)  \right]  $
\begin{align}
\mathcal{E}\left(  \mathbf{n},\mathbf{\psi},\mathbf{\theta}\right)   &
=\frac{\mathcal{H}\left(  \mathbf{n},\mathbf{\psi},\mathbf{\theta}\right)
}{S_{N}\left(  \mathbf{n}\right)  }=\frac{1}{S_{N}\left(  \mathbf{n}\right)
}\left\{  \sum_{1\leq j\leq s}n_{j}S_{N-1}\left(  \jmath\right)  \left\{
\overset{}{\left\langle \psi_{j}\left\vert h_{1}\psi_{j}\right.  \right\rangle
+\left\langle \psi_{j+s}\left\vert h_{1}\psi_{j+s}\right.  \right\rangle
}\right.  \right.  \nonumber\\
&  \left.  +\overset{}{\,\left\langle \psi_{j}\psi_{j+s}|h_{2}|\psi_{j}%
\psi_{j+s}\right\rangle }\right\}  +\sum_{1\leq j,k\leq s}n_{j}^{\frac{1}{2}%
}n_{k}^{\frac{1}{2}}e^{i\left(  \theta_{j}-\theta_{k}\right)  }S_{N-1}\left(
\jmath k\right)  \left\langle \psi_{j}\psi_{j+s}|h_{2}|\psi_{k}\psi
_{k+s}\right\rangle \nonumber\\
&  \qquad\qquad\qquad\qquad\qquad\left.  +\sum_{1\leq j<k\leq s}n_{i}%
n_{j}S_{N-2}\left(  \jmath k\right)  \left\langle \psi_{j}\psi_{k}|||\psi
_{j}\psi_{k}\right\rangle \right\}  ,\label{thetaen}%
\end{align}
where $\left\vert \mathbf{\psi}\right\rangle $ is a representative set of
CGSO's for the class $\left[  D^{1}\left(  g^{N}\left(  \mathbf{0}\right)
\right)  \right]  $. As $\left(  \mathbf{n},\mathbf{\psi}\right)  $ uniquely
defines a $2N$-electron AGP-FORDO\ the expression eq. $\left(  \ref{thetaen}%
\right)  $ can be used to define a density matrix energy functional
\begin{equation}
\mathcal{F}\left(  D^{1}\right)  =\min_{\mathbf{\theta\in}\Theta}%
\mathcal{E}\left(  \mathbf{n},\mathbf{\psi},\mathbf{\theta}\right)
\end{equation}
over the manifold of AGP\ states, where
\begin{equation}
\Theta=\left\{  \left.  \theta_{j}\right\vert \theta_{1}=0,\,0<\theta_{j}%
\leq2\pi,2\leq j\leq s\right\}  .
\end{equation}
Noting that each class is determined by $\left(  \mathbf{n},\mathbf{\psi
}\right)  $ this functional can then be minimized to give
\begin{equation}
\min_{D^{1}\in\mathcal{S}_{1}}\mathcal{F}\left(  D^{1}\right)  =\min
_{\mathbf{n\in}\mathcal{D}_{n},\left\vert \mathbf{\psi}\right\rangle
\in\mathcal{C}}\mathcal{F}\left(  D^{1}\left(  \mathbf{n},\mathbf{\psi
}\right)  \right)  =\min_{\left\vert g\right\rangle \in\mathcal{H}^{2}%
}\mathcal{E}\left(  \left\vert g^{N}\right\rangle \right)  ,\label{MinDMF}%
\end{equation}
where the domains $\mathcal{D}_{n}$, $\mathcal{C}$ and $\mathcal{S}_{1}$ are
given as%
\begin{align}
\mathcal{D}_{n} &  =\left\{  \left.  n_{j}\right\vert ;\sum_{1\leq j\leq
s}n_{j}=1,n_{j}\geq0,\quad1\leq j\leq s\right\}  ,\\
\mathcal{C} &  =\left\{  \left.  \left\vert \mathbf{\psi}\right\rangle
\right\vert ;\left\vert \mathbf{\psi}\right\rangle \neq e^{i\mathbf{\theta}%
}\left\vert \mathbf{\psi}^{\prime}\right\rangle ,\text{ when }\left\vert
\mathbf{\psi}^{\prime}\right\rangle \in\mathcal{C}\right\}
\end{align}
and
\begin{equation}
\mathcal{S}_{1}=\left\{  \left.  D^{1}\right\vert ;Tr\left\{  D^{1}\right\}
=2N,D^{1}\geq0,D^{1}\text{ at least doubly degenerate}\right\}  .
\end{equation}
The solution to eq. $\left(  \ref{MinDMF}\right)  $ characterizes the
$2N$-electron AGP state that minimizes the total energy.

\subsection{Energy Functional in Terms of Group Parameters}

The manifold of AGP\ states can be parameterized 

\section{Necessary Conditions for the Minimization of the Energy
\label{Evariation}}

\subsection{Variational Problem}

The solution to the constrained optimization problem
\begin{equation}
\min_{\left(  \mathbf{n,\psi,\theta}\right)  }\mathcal{E}\left(
\mathbf{n,\psi,\theta}\right)  \label{min}%
\end{equation}
subject to the phase constraints%
\begin{equation}
\theta_{1}=0;\,0<\theta_{j}\leq2\pi,2\leq j\leq s, \label{thetacon}%
\end{equation}
the orthonormality constraints%
\begin{equation}
\left\langle \left.  \psi_{j}\right\vert \psi_{k}\right\rangle =\delta
_{jk};\quad1\leq j,k\leq2s, \label{orthonorm}%
\end{equation}
phase inequivalent constraint
\begin{equation}
\left\vert \mathbf{\psi}\right\rangle \mathbf{\in\,}\mathcal{C} \label{rep}%
\end{equation}
and the positivity and normalization geminal occupation constraint
\begin{subequations}
\begin{align}
n_{j}  &  \geq0;\quad1\leq j\leq s,\label{pos1}\\
\sum_{1\leq j\leq s}n_{j}  &  =1 \label{gemnorm}%
\end{align}
characterize an optimal $2N$-electron AGP\ state.

Some comments on the minimization problem eq. $\left(  \ref{min}\right)  $:
\end{subequations}
\begin{enumerate}
\item The number of electrons being described is implicit in the definition of
$\mathcal{E}\left(  \mathbf{n,\psi,\theta}\right)  $ and one could show this
dependency explicitly as
\begin{equation}
\mathcal{E}=\mathcal{E}\left(  2N,\mathbf{n,\psi,\theta}\right)  .
\end{equation}
This is useful if one needs to consider an explicit form for the chemical
potential $\frac{\partial\mathcal{E}}{\partial N}.$

\item The constraint eq. $\left(  \ref{thetacon}\right)  $ is always satisfied
as $\mathcal{E}$ is a periodic function of $\left\{  \left.  \theta
_{j}\right\vert 2\leq j\leq s\right\}  $ with period $2\pi$ in each variable.

\item We show in section \ref{Fock} that the normalization constraints
\begin{equation}
\left\langle \left.  \psi_{j}\right\vert \psi_{j}\right\rangle =1;\quad1\leq
j\leq2s \label{norm}%
\end{equation}
imply the orthonormality constraints
\begin{equation}
\left\langle \left.  \psi_{j}\right\vert \psi_{k}\right\rangle =\delta
_{jk};\quad1\leq j,k\leq2s.
\end{equation}


\item In the optimization procedure we describe we allow the constraint eq.
$\left(  \ref{rep}\right)  $ to relax, however we show that this does effect
the final solution.
\end{enumerate}

In section \ref{DMF} we obtain an optimal value, $\mathbf{\theta}_{\#}%
\in\Theta$ , of $\mathbf{\theta}$ for fixed values of $\left(  \mathbf{n,\psi
}\right)  $ by numerically solving the stationarity equation for
$\mathbf{\theta}$ and in section \ref{gemocc} we utilize and describe an
iterative eigenvector procedure to determine an optimal value $\mathbf{n}%
_{\#}$, subject to the constraints eqs. $\left(  \ref{pos1}\right)  $ and
$\left(  \ref{gemnorm}\right)  $ for fixed values of $\left(  \mathbf{\psi
,\theta}\right)  $. In section \ref{Fock} the we utilize the method of
Lagrange multipliers based on the Lagrangian
\begin{equation}
\mathcal{\tilde{E}}\left(  \mathbf{n,\psi,\theta}\right)  =\mathcal{E}\left(
\mathbf{n,\psi,\theta}\right)  -\sum_{1\leq j\leq2s}\varepsilon_{j}\left(
\left\langle \left.  \psi_{j}\right\vert \psi_{j}\right\rangle -1\right)
\label{larE}%
\end{equation}
to obtain an optimal set $\left\{  \left.  \psi_{\#j}\right\vert 1\leq
j\leq2s\right\}  $ of orthonormal CGSO's for fixed values of $\left(
\mathbf{\psi,\theta}\right)  $.

\subsection{Stationarity Conditions}

A necessary condition for $\left(  \mathbf{n}_{\#}\mathbf{,\psi}%
_{\#}\mathbf{,\theta}_{\#}\right)  $ to be a solution of eq. $\left(
\ref{min}\right)  $ is that there exists Lagrange multipliers $\left\{
\left.  \varepsilon_{j}\right\vert 1\leq j\leq2s\right\}  $ such that
\begin{equation}
\partial^{1}\mathcal{\tilde{E}}\left(  \mathbf{n}_{\#}\mathbf{,\psi}%
_{\#}\mathbf{,\theta}_{\#}\right)  =\mathbf{0,}%
\end{equation}
where
\begin{align}
&  \left.  \partial^{1}\mathcal{\tilde{E}}\left(  \mathbf{n}_{\#}%
\mathbf{,\psi}_{\#}\mathbf{,\theta}_{\#}\right)  \right.  =\nonumber\\
&  \left(  \partial_{\mathbf{n}}^{1}\mathcal{\tilde{E}}\left(  \mathbf{n}%
_{\#}\mathbf{,\psi}_{\#}\mathbf{,\theta}_{\#}\right)  ,\partial_{\mathbf{\psi
}}^{1}\mathcal{\tilde{E}}\left(  \mathbf{n}_{\#}\mathbf{,\psi}_{\#}%
\mathbf{,\theta}_{\#}\right)  ,\partial_{\mathbf{\theta}}^{1}\mathcal{\tilde
{E}}\left(  \mathbf{n}_{\#}\mathbf{,\psi}_{\#}\mathbf{,\theta}_{\#}\right)
\right)  ,
\end{align}
and
\begin{align}
\partial_{\mathbf{n}}^{1}\mathcal{\tilde{E}}\left(  \mathbf{n}_{\#}%
\mathbf{,\psi}_{\#}\mathbf{,\theta}_{\#}\right)  _{j}  &  =\frac
{\partial\mathcal{\tilde{E}}\left(  \mathbf{n}_{\#}\mathbf{,\psi}%
_{\#}\mathbf{,\theta}_{\#}\right)  }{\partial n_{j}};\quad1\leq j\leq
s\nonumber\\
\partial_{\mathbf{\psi}}^{1}\mathcal{\tilde{E}}\left(  \mathbf{n}%
_{\#}\mathbf{,\psi}_{\#}\mathbf{,\theta}_{\#}\right)  _{j}  &  =\frac
{\delta\mathcal{\tilde{E}}\left(  \mathbf{n}_{\#}\mathbf{,\psi}_{\#}%
\mathbf{,\theta}_{\#}\right)  }{\delta\psi_{j}^{\ast}};\quad1\leq j\leq
s\nonumber\\
\partial_{\mathbf{\psi}}^{1}\mathcal{\tilde{E}}\left(  \mathbf{n}%
_{\#}\mathbf{,\psi}_{\#}\mathbf{,\theta}_{\#}\right)  _{j+s}  &  =\frac
{\delta\mathcal{\tilde{E}}\left(  \mathbf{n}_{\#}\mathbf{,\psi}_{\#}%
\mathbf{,\theta}_{\#}\right)  }{\delta\psi_{j}};\quad1\leq j\leq s\nonumber\\
\partial_{\mathbf{\theta}}^{1}\mathcal{\tilde{E}}\left(  \mathbf{n}%
_{\#}\mathbf{,\psi}_{\#}\mathbf{,\theta}_{\#}\right)  _{j}  &  =\frac
{\partial\mathcal{\tilde{E}}\left(  \mathbf{n}\left(  \mathbf{\alpha}%
_{\#}\right)  \mathbf{,\psi}_{\#}\mathbf{,\theta}_{\#}\right)  }%
{\partial\theta_{j}};\quad2\leq j\leq s
\end{align}
and that $\left(  \mathbf{n}_{\#}\mathbf{,\psi}_{\#}\mathbf{,\theta}%
_{\#}\right)  $ satisfy the constraints eqs. $\left(  \ref{thetacon}\right)
,\left(  \ref{pos1}\right)  ,\left(  \ref{gemnorm}\right)  $ and $\left(
\ref{orthonorm}\right)  $.

As is well known this condition only characterizes constrained stationary
points of $\mathcal{E}$ and the Hessian matrix of second derivatives has to be
examined in order to determine if the point $\left(  \mathbf{n}_{\#}%
\mathbf{,\psi}_{\#}\mathbf{,\theta}_{\#}\right)  $ does in fact minimize
$\mathcal{E}.$

\section{Generalized Fock Operator \label{Fock} for the Variation wrt
$\mathbf{\psi}$}

\subsection{Functional Derivatives}

The energy expression eq. $\left(  \ref{energy}\right)  $ can be
differentiated wrt $\mathbf{\psi}^{\ast},$ $\mathbf{\psi}$ to produce the
functional derivatives (Appendix )
\begin{equation}
\left\{  \frac{\delta\mathcal{E}\left(  \mathbf{n,\psi,\theta}\right)
}{\delta\psi_{j}^{\ast}},\frac{\delta\mathcal{E}\left(  \mathbf{n,\psi,\theta
}\right)  }{\delta\psi_{j}};\quad1\leq j\leq s\right\}  ,
\end{equation}
which can be expressed as%
\begin{align}
\frac{\delta\mathcal{E}\left(  \mathbf{n,\psi,\theta}\right)  }{\delta\psi
_{j}^{\ast}}  &  =\frac{1}{S_{N}\left(  \mathbf{n}\right)  }\frac
{\delta\mathcal{H}\left(  \mathbf{n,\psi,\theta}\right)  }{\delta\psi
_{j}^{\ast}}\\
\frac{\delta\mathcal{E}\left(  \mathbf{n,\psi,\theta}\right)  }{\delta\psi
_{j}}  &  =\frac{1}{S_{N}\left(  \mathbf{n}\right)  }\frac{\delta
\mathcal{H}\left(  \mathbf{n,\psi,\theta}\right)  }{\delta\psi_{j}}%
\end{align}
in terms of the derivatives of the unnormalized energy $\mathcal{H}\left(
\mathbf{n,\psi,\theta}\right)  ,$ as $S_{N}\left(  \mathbf{n}\right)  $ does
not depend on $\left(  \mathbf{\psi,\theta}\right)  $ when $\left\{  \left.
\left\vert \psi_{j}\right\rangle \right\vert \quad1\leq j\leq2s\right\}  $ are
orthonormal. It is easy to show that
\begin{equation}
\frac{\delta\mathcal{H}\left(  \mathbf{n,\psi,\theta}\right)  }{\delta\psi
_{j}^{\ast}}=\left(  \frac{\delta\mathcal{H}\left(  \mathbf{n,\psi,\theta
}\right)  }{\delta\psi_{j}}\right)  ^{\ast}%
\end{equation}
so we only need consider $\left\{  \frac{\delta\mathcal{H}\left(
\mathbf{n,\psi,\theta}\right)  }{\delta\psi_{j}^{\ast}}\right\}  .$ These
derivatives can be expressed $\left(  \text{see Appendix \ref{variational}%
}\right)  $ in terms of a generalized self adjoint Fock operator $F\left(
\mathbf{n,\psi,\theta}\right)  $ as
\begin{equation}
\frac{\delta\mathcal{H}\left(  \mathbf{n,\psi,\theta}\right)  }{\delta\psi
_{j}^{\ast}}=F\left(  \mathbf{n,\psi,\theta}\right)  \left\vert \psi
_{j}\right\rangle ,\quad1\leq j\leq2s,
\end{equation}
that depends on the occupation numbers $\mathbf{n,}$ representative CGSO's
$\mathbf{\psi,}$ of the class $\left[  D_{2N}^{1}\right]  $ and the internal
phases $\mathbf{\theta}$ and is given by
\begin{equation}
F\left(  \mathbf{n,\psi,\theta}\right)  =\sum_{1\leq j,m\leq2s}\left\{
n_{j}S_{N-1}\left(  \jmath\right)  h_{1mj}+A_{mj}-B_{mj}+\left\{
C_{mj}-D_{mj}\right\}  +2\overset{}{\left\{  E_{mj}-F_{mj}\right\}  }\right\}
\left\vert \psi_{m}\right\rangle \left\langle \psi_{j}\right\vert .
\label{genfock}%
\end{equation}
The matrices $\left\{  \mathbb{A},\mathbb{B},\mathbb{C},\mathbb{D}%
,\mathbb{E},\mathbb{F}\right\}  $ are given in terms of the matrix elements
$\left\{  \left.  T_{jklm}\right\vert 1\leq j,k,l,m\leq2s\right\}  ,$ which
are defined as
\begin{equation}
T_{jklm}=\left(  \left.  \psi_{j}\psi_{k}\right\vert h_{2}\psi_{l}\psi
_{m}\right)  =\int\psi_{l}^{\ast}\left(  \mathbf{x}_{1}\right)  \psi_{j}%
^{\ast}\left(  \mathbf{x}_{2}\right)  \frac{1}{\left\Vert \mathbf{r}%
_{1}-\mathbf{r}_{2}\right\Vert }\psi_{m}\left(  \mathbf{x}_{1}\right)
\psi_{k}\left(  \mathbf{x}_{2}\right)  d\mathbf{x}_{1}d\mathbf{x}_{2}.
\end{equation}
The matrix elements of $\left\{  \mathbb{A},\mathbb{B},\mathbb{C}%
,\mathbb{D},\mathbb{E},\mathbb{F}\right\}  $ for $1\leq m\leq2s$ are given by
\begin{equation}
A_{mj}=\left\{
\begin{array}
[c]{c}%
n_{j}S_{N-1}\left(  \jmath\right)  T_{mj+sjj+s};\,1\leq j\leq s\\
n_{j}S_{N-1}\left(  \jmath\right)  T_{mj-sjj-s};\,s+1\leq j\leq2s
\end{array}
\right.  , \label{a1}%
\end{equation}%
\begin{equation}
B_{mj}=\left\{
\begin{array}
[c]{c}%
n_{j}S_{N-1}\left(  \jmath\right)  T_{mj+sj+sj};\,1\leq j\leq s\\
n_{j}S_{N-1}\left(  \jmath\right)  T_{mj-sj-sj};\,s+1\leq j\leq2s
\end{array}
\right.  , \label{b1}%
\end{equation}
noting that $n_{j+s}=n_{j}$ and $\theta_{j+s}=\theta_{j}$ for $1\leq j\leq s$
we have the phase dependent $2$-electron terms
\begin{equation}
C_{mj}=\left\{
\begin{array}
[c]{c}%
%TCIMACRO{\tsum \limits_{1\leq k\leq s}}%
%BeginExpansion
{\textstyle\sum\limits_{1\leq k\leq s}}
%EndExpansion
n_{j}^{\frac{1}{2}}n_{k}^{\frac{1}{2}}S_{N-1}\left(  \jmath k\right)
e^{i\left(  \theta_{j}-\theta_{k}\right)  }T_{mj+skk+s};\,1\leq j\leq s\\%
%TCIMACRO{\tsum \limits_{1\leq k\leq s}}%
%BeginExpansion
{\textstyle\sum\limits_{1\leq k\leq s}}
%EndExpansion
n_{j}^{\frac{1}{2}}n_{k}^{\frac{1}{2}}S_{N-1}\left(  \jmath k\right)
e^{i\left(  \theta_{j}-\theta_{k}\right)  }T_{mj-sk+sk};\,s+1\leq j\leq2s
\end{array}
\right.  , \label{c1}%
\end{equation}%
\begin{equation}
D_{mj}=\left\{
\begin{array}
[c]{c}%
%TCIMACRO{\tsum \limits_{1\leq k\leq s}}%
%BeginExpansion
{\textstyle\sum\limits_{1\leq k\leq s}}
%EndExpansion
n_{j}^{\frac{1}{2}}n_{k}^{\frac{1}{2}}S_{N-1}\left(  \jmath k\right)
e^{i\left(  \theta_{j}-\theta_{k}\right)  }T_{mj+sk+sk};\,1\leq j\leq s\\%
%TCIMACRO{\tsum \limits_{1\leq k\leq s}}%
%BeginExpansion
{\textstyle\sum\limits_{1\leq k\leq s}}
%EndExpansion
n_{j}^{\frac{1}{2}}n_{k}^{\frac{1}{2}}S_{N-1}\left(  \jmath k\right)
e^{i\left(  \theta_{j}-\theta_{k}\right)  }T_{mj-skk+s};\,s+1\leq j\leq2s
\end{array}
\right.  , \label{d1}%
\end{equation}
the phase independent $2$-electron terms
\begin{equation}
F_{mj}=\sum_{1\leq k\leq2s}n_{j}n_{k}S_{N-2}\left(  \jmath k\right)
T_{mkkj};\,1\leq j\leq2s, \label{f1}%
\end{equation}
and%
\begin{equation}
E_{mj}=\sum_{1\leq k\leq2s}n_{j}n_{k}S_{N-2}\left(  \jmath k\right)
T_{mkjk};\,1\leq j\leq2s. \label{e1}%
\end{equation}


\subsection{Orthonormality Constraints}

A necessary condition for $\mathbf{\psi}_{\#}$ $=\mathbf{\psi}_{\#}\left(
\mathbf{n,\theta}\right)  $ to be a orthonormally constrained stationary point
of $\mathcal{E}\left(  \mathbf{n,\psi,\theta}\right)  $ for a fixed value of
$\left(  \mathbf{n,\theta}\right)  $ is given by
\begin{subequations}
\begin{align}
\partial_{\mathbf{\psi}}^{1}\mathcal{\tilde{H}}\left(  \mathbf{n,\mathbf{\psi
}_{\#},\theta}\right)   &  =\mathbf{0}\label{nec1}\\
\left\langle \left.  \psi_{\#j}\right\vert \psi_{\#k}\right\rangle  &
=\delta_{jk};1\leq j,k\leq2s,
\end{align}
where the Lagrangian $\mathcal{\tilde{H}}\left(  \mathbf{n,\mathbf{\psi
},\theta}\right)  $ is defined as
\end{subequations}
\begin{equation}
\mathcal{\tilde{H}}\left(  \mathbf{n,\mathbf{\psi},\theta}\right)
=\mathcal{H}\left(  \mathbf{n,\mathbf{\psi},\theta}\right)  -\sum_{1\leq
j\leq2s}\varepsilon_{j}\left\{  \left\langle \left.  \psi_{j}\right\vert
\psi_{j}\right\rangle -1\right\}  , \label{lagH}%
\end{equation}
as
\begin{equation}
\partial_{\mathbf{\psi}}^{1}\mathcal{\tilde{H}}\left(  \mathbf{n,\mathbf{\psi
},\theta}\right)  =\partial_{\mathbf{\psi}}^{1}\mathcal{\tilde{E}}\left(
\mathbf{n,\mathbf{\psi},\theta}\right)
\end{equation}
for orthonormally constrained $\mathbf{\mathbf{\psi.}}$ The condition eq.
$\left(  \ref{nec1}\right)  $ can be expressed in terms of the generalized
Fock operator as
\begin{equation}
F\left(  \mathbf{n,\mathbf{\psi},\theta}\right)  \left\vert \psi
_{j}\right\rangle =\varepsilon_{j}\left\vert \psi_{j}\right\rangle ;1\leq
j\leq2s. \label{eig}%
\end{equation}
As $F\left(  \mathbf{\psi,n}\right)  =F\left(  \mathbf{\psi,n}\right)
^{\dagger}$ this implies that one can find solutions $\left\{  \left.
\left\vert \psi_{m}\right\rangle \right\vert 1\leq m\leq2s\right\}  $ for
$\left\{  \left.  \varepsilon_{m}\right\vert 1\leq m\leq2s\right\}  $ such
that
\begin{equation}
\left\langle \left.  \psi_{j}\right\vert \psi_{k}\right\rangle =\delta
_{jk},;\quad1\leq j,k\leq2s.
\end{equation}
Thus one can always find stationary points that satisfy the \emph{normality}
constraints
\begin{equation}
\left\langle \left.  \psi_{j}\right\vert \psi_{j}\right\rangle =1;\quad1\leq
j\leq2s
\end{equation}
that also satisfy the orthonormality constraints eq. $\left(  \ref{norm}%
\right)  .$ This conclusion is also consistent with the observation that the
orthogonality constraints are \emph{not }active for $\left\{  \left.
\left\vert \psi_{j}\right\rangle \right\vert \right\}  $ that satisfy eq.
$\left(  \ref{eig}\right)  $ as "off diagonal" Lagrange multipliers,$\left\{
\left.  \varepsilon_{jk}\right\vert j\neq k\right\}  $ that could have been
associated with such constraints would have the value zero.

The Lagrangian eq. $\left(  \ref{lagH}\right)  $ can also be used in the
IPS\ case to produce the standard Hartree-Fock equations for $2N$-electrons
instead of the traditional Lagrangian
\begin{equation}
\mathcal{\breve{H}}\left(  \mathbf{n}_{IP}\mathbf{,\mathbf{\psi},0}\right)
=\mathcal{H}\left(  \mathbf{n}_{IP}\mathbf{,\mathbf{\psi},0}\right)
-\sum_{1\leq j,k\leq2N}\varepsilon_{jk}\left\{  \left\langle \left.  \psi
_{j}\right\vert \psi_{k}\right\rangle -\delta_{jk}\right\}  , \label{tradlag}%
\end{equation}
that explicitly holds the orthonormality constraints between occupied
orbitals. \emph{However }using a Lagrangian that explicitly holds the
orthonormality constraints in the AGP\ case \emph{does not} lead to equations
that impose conditions on $\mathbf{\mathbf{\psi}}$.

\section{Density Matrix Functional\label{DMF}}

The AGP states that have the same FORDO are parameterized by a vector
$\mathbf{\theta}$ and have energies given by
\begin{equation}
\tilde{f}\left(  \mathbf{\theta}\right)  =K\left(  \mathbf{n},\mathbf{\psi
}\right)  +\frac{1}{S_{N}\left(  \mathbf{n}\right)  }\sum_{1\leq j,k\leq
s}n_{j}^{\frac{1}{2}}n_{k}^{\frac{1}{2}}e^{i\left(  \theta_{j}-\theta
_{k}\right)  }S_{N-1}\left(  \jmath k\right)  \left\langle \psi_{j}\psi
_{j+s}|h_{2}\psi_{k}\psi_{k+s}\right\rangle ,
\end{equation}
where the constant $K$ is given by eq. $\left(  \ref{thetaen}\right)  $ and
\begin{equation}
\mathbf{\theta}=\left(  0,\theta_{2},\cdots,\theta_{s}\right)  ;\quad
0<\theta_{j}\leq2\pi,\quad2\leq j\leq s.
\end{equation}
As $K\left(  \mathbf{n},\mathbf{\psi}\right)  $ is constant for fixed $\left(
\mathbf{n},\mathbf{\psi}\right)  $ we are interested in the minimum of
\begin{align}
f\left(  \mathbf{\theta}\right)   &  =\frac{1}{S_{N}\left(  \mathbf{n}\right)
}\sum_{1\leq j,k\leq s}n_{j}^{\frac{1}{2}}n_{k}^{\frac{1}{2}}e^{i\left(
\theta_{j}-\theta_{k}\right)  }S_{N-1}\left(  \jmath k\right)  \left\langle
\psi_{j}\psi_{j+s}\right\vert \left\vert \psi_{k}\psi_{k+s}\right\rangle
\nonumber\\
&  =\sum_{1\leq j,k\leq s}e^{-i\theta_{k}}M_{kj}e^{i\theta_{j}}=\left(
e^{i\mathbf{\theta}}\right)  ^{\dagger}\mathbb{M}e^{i\mathbf{\theta}},
\end{align}
where $e^{i\mathbf{\theta}}$ is the vector
\begin{align}
\left(  e^{i\mathbf{\theta}}\right)  _{j}  &  =e^{i\theta_{j}};\quad2\leq
j\leq s\nonumber\\
\left(  e^{i\mathbf{\theta}}\right)  _{1}  &  =1
\end{align}
and $\mathbb{M}$ the positive matrix given by
\begin{equation}
M_{kj}=\frac{1}{S_{N}\left(  \mathbf{n}\right)  }n_{j}^{\frac{1}{2}}%
n_{k}^{\frac{1}{2}}S_{N-1}\left(  \jmath k\right)  \left\langle \psi_{j}%
\psi_{j+s}\right\vert \left\vert \psi_{k}\psi_{k+s}\right\rangle .
\end{equation}
We can expand $f$ in a Taylors series about $\mathbf{\theta}_{0}$ to obtain
\begin{equation}
f\left(  \mathbf{\theta}_{0}+\delta\mathbf{\theta}\right)  =f\left(
\mathbf{\theta}_{0}\right)  +\partial_{\mathbf{\theta}}f\left(  \mathbf{\theta
}_{0}\right)  ^{t}\delta\mathbf{\theta}+\frac{1}{2}\delta\mathbf{\theta}%
^{t}\partial_{\mathbf{\theta}}^{2}f\left(  \mathbf{\theta}_{0}\right)
\delta\mathbf{\theta+\cdots,}%
\end{equation}
where
\begin{equation}
\partial_{\mathbf{\theta}}f\left(  \mathbf{\theta}_{0}\right)  _{l}%
=\frac{\partial f\left(  \mathbf{\theta}_{0}\right)  }{\partial\theta_{l}%
}=i\sum_{1\leq j\leq s}\left\{  e^{-i\theta_{0j}}M_{jl}e^{i\theta_{0l}%
}-e^{-i\theta_{0l}}M_{lj}e^{i\theta_{0j}}\right\}  =-2\sum_{1\leq j\leq
s}\operatorname{Im}\left\{  e^{-i\theta_{0j}}M_{jl}e^{i\theta_{0l}}\right\}
\end{equation}
and
\begin{align}
\partial_{\mathbf{\theta}}^{2}f\left(  \mathbf{\theta}_{0}\right)  _{lm}  &
=\frac{\partial^{2}f\left(  \mathbf{\theta}_{0}\right)  }{\partial\theta
_{l}\partial\theta_{m}}\nonumber\\
&  =2\left\{  \operatorname{Im}\left\{  \sum_{1\leq j\leq s}e^{-i\theta_{0j}%
}M_{jl}e^{i\theta_{0l}}\right\}  \delta_{lm}-\operatorname{Im}\left\{
e^{-i\theta_{0l}}M_{lm}e^{i\theta_{0m}}\right\}  \right\}  .
\end{align}
The stationary points $\mathbf{\theta}_{\#}$ of $f$ are given by
\begin{equation}
\partial_{\mathbf{\theta}}f\left(  \mathbf{\theta}_{\#}\right)  =\mathbf{0}
\label{gradf}%
\end{equation}
and these are local minimum if
\begin{equation}
\partial_{\mathbf{\theta}}^{2}f\left(  \mathbf{\theta}_{\#}\right)
\geq\mathbf{0.}%
\end{equation}
The eq. $\left(  \ref{gradf}\right)  $ can be solved numerically by any
standard optimization technique, though it would be more convenient if an
analytic solution to eq. $\left(  \ref{gradf}\right)  $ could be found.

\section{Variation of the Occupation Numbers\label{OccVar}}

In this section we show how by optimizing wrt geminal coefficients $\left\{
c_{j}\right\}  $ for a fixed value of $\left(  \mathbf{\mathbf{\psi},\theta
}\right)  $ one can determine an optimal value of $\mathbf{n}$ that satisfies
the positivity and normalization constraints.

The $2N$-electron, unnormalized, AGP state, $\left\vert g^{N}\right\rangle $,
can be expressed as
\begin{equation}
\left\vert g^{N}\right\rangle =\left(  G^{\dagger}\right)  ^{N}\left\vert
\phi\right\rangle ,
\end{equation}
where
\begin{equation}
G^{\dagger}=\sum_{1\leq j\leq s}c_{j}b_{j}^{\dagger}b_{j+s}^{\dagger},
\label{gcreat}%
\end{equation}
and the electron creation operators $\left\{  \left.  b_{j}^{\dagger
}\right\vert 1\leq j\leq2s\right\}  $ produce the CGSO's $\left\{  \left.
\left\vert \psi_{j}\right\rangle \right\vert 1\leq j\leq2s\right\}  $ by
acting on the zero particle vacuum vector i.e.
\begin{equation}
\left\vert \psi_{j}\right\rangle =b_{j}^{\dagger}\left\vert \phi\right\rangle
;\quad1\leq j\leq2s.
\end{equation}
For fixed $\left(  \mathbf{\mathbf{\psi},\theta}\right)  $ the geminal creator
in eq. $\left(  \ref{gcreat}\right)  $ is a function of the canonical
coefficients $\left\{  \left.  c_{j}\right\vert 1\leq j\leq s\right\}  ,$
which can be differentiated to give%
\begin{equation}
\frac{\partial G\left(  \mathbf{c}\right)  ^{\dagger N}}{\partial c_{j}%
}=Nb_{j}^{\dagger}b_{j+s}^{\dagger}G\left(  \mathbf{c}\right)  ^{\dagger N-1},
\end{equation}
which leads to
\begin{equation}
G\left(  \mathbf{c}\right)  ^{\dagger N}=\frac{1}{N}\sum_{1\leq j\leq s}%
c_{j}\frac{\partial G\left(  \mathbf{c}\right)  ^{\dagger N}}{\partial c_{j}}.
\label{euler}%
\end{equation}
The energy expression
\begin{equation}
\mathcal{E}\left(  \mathbf{c}\right)  =\frac{\mathcal{H}\left(  \mathbf{c}%
\right)  }{S_{N}\left(  \mathbf{c}\right)  }=\frac{\left\langle G\left(
\mathbf{c}\right)  ^{\dagger N}\phi|HG\left(  \mathbf{c}\right)  ^{\dagger
N}\phi\right\rangle }{\left\langle G\left(  \mathbf{c}\right)  ^{\dagger
N}\phi|G\left(  \mathbf{c}\right)  ^{\dagger N}\phi\right\rangle }%
\end{equation}
produces the derivative
\begin{equation}
\partial_{\mathbf{c}}^{1}\mathcal{E}\left(  \mathbf{c}\right)  =\frac{1}%
{S_{N}\left(  \mathbf{c}\right)  }\left\{  \partial_{\mathbf{c}}%
^{1}\mathcal{H}\left(  \mathbf{c}\right)  -\mathcal{E}\left(  \mathbf{c}%
\right)  \partial_{\mathbf{c}}^{1}S_{N}\left(  \mathbf{c}\right)  \right\}  ,
\end{equation}
where
\begin{equation}
\partial_{c_{j}}^{1}\mathcal{H}\left(  \mathbf{c}\right)  =N\left\{
\left\langle b_{j}^{\dagger}b_{j+s}^{\dagger}G\left(  \mathbf{c}\right)
^{\dagger N-1}\phi|HG\left(  \mathbf{c}\right)  ^{\dagger N}\phi\right\rangle
+\left\langle G\left(  \mathbf{c}\right)  ^{\dagger N}\phi|Hb_{j}^{\dagger
}b_{j+s}^{\dagger}G\left(  \mathbf{c}\right)  ^{\dagger N-1}\phi\right\rangle
\right\}
\end{equation}
and
\begin{equation}
\partial_{c_{j}}^{1}S_{N}\left(  \mathbf{c}\right)  =N\left\{  \left\langle
b_{j}^{\dagger}b_{j+s}^{\dagger}G\left(  \mathbf{c}\right)  ^{\dagger N-1}%
\phi|G\left(  \mathbf{c}\right)  ^{\dagger N}\phi\right\rangle +\left\langle
G\left(  \mathbf{c}\right)  ^{\dagger N}\phi|b_{j}^{\dagger}b_{j+s}^{\dagger
}G\left(  \mathbf{c}\right)  ^{\dagger N-1}\phi\right\rangle \right\}  .
\end{equation}
Using eq. $\left(  \ref{euler}\right)  $ one obtains
\begin{equation}
\frac{1}{2N}\sum_{1\leq j\leq s}c_{j}\frac{\partial\mathcal{E}\left(
\mathbf{c}\right)  }{\partial c_{j}}=\mathbf{c}^{t}\partial_{\mathbf{c}}%
^{1}\mathcal{E}\left(  \mathbf{c}\right)  =\frac{1}{2NS_{N}\left(
\mathbf{c}\right)  }\left\{  \mathcal{H}\left(  \mathbf{c}\right)
-\mathcal{E}\left(  \mathbf{c}\right)  S_{N}\left(  \mathbf{c}\right)
\right\}  =0,
\end{equation}
i.e. the vector $\mathbf{c}$ is always perpendicular to the vector
$\partial_{\mathbf{c}}^{1}\mathcal{E}\left(  \mathbf{c}\right)  $.

The second derivatives $\left\{  \partial_{c_{j}c_{k}}^{2}\mathcal{E}\left(
\mathbf{c}\right)  \right\}  $ are given by
\begin{equation}
\partial_{\mathbf{c}}^{2}\mathcal{E}\left(  \mathbf{c}\right)  =\frac{1}%
{S_{N}\left(  \mathbf{c}\right)  }\left\{  \partial_{\mathbf{c}}%
^{2}\mathcal{H}\left(  \mathbf{c}\right)  -\mathcal{E}\left(  \mathbf{c}%
\right)  \partial_{\mathbf{c}}^{2}S_{N}\left(  \mathbf{c}\right)
-\partial_{\mathbf{c}}^{1}\mathcal{E}\left(  \mathbf{c}\right)  \otimes
\partial_{\mathbf{c}}^{1}S_{N}\left(  \mathbf{c}\right)  -\partial
_{\mathbf{c}}^{1}S_{N}\left(  \mathbf{c}\right)  \otimes\partial_{\mathbf{c}%
}^{1}\mathcal{E}\left(  \mathbf{c}\right)  \right\}  , \label{Hess}%
\end{equation}
where
\begin{align}
\partial_{c_{j}c_{k}}^{2}\mathcal{H}\left(  \mathbf{c}\right)   &
=N^{2}\left\{  \left\langle b_{j}^{\dagger}b_{j+s}^{\dagger}G\left(
\mathbf{c}\right)  ^{\dagger N-1}\phi|Hb_{k}^{\dagger}b_{k+s}^{\dagger
}G\left(  \mathbf{c}\right)  ^{\dagger N-1}\right\rangle +\left\langle
b_{k}^{\dagger}b_{k+s}^{\dagger}G\left(  \mathbf{c}\right)  ^{\dagger N-1}%
\phi|Hb_{j}^{\dagger}b_{j+s}^{\dagger}G\left(  \mathbf{c}\right)  ^{\dagger
N-1}\right\rangle \right\} \nonumber\\
&  +N\left(  N-1\right)  \left\{  \left\langle G\left(  \mathbf{c}\right)
^{\dagger N}\phi|Hb_{j}^{\dagger}b_{j+s}^{\dagger}b_{k}^{\dagger}%
b_{k+s}^{\dagger}G\left(  \mathbf{c}\right)  ^{\dagger N-2}\phi\right\rangle
\right. \nonumber\\
&  \qquad\qquad\qquad\qquad\qquad\qquad\qquad\qquad\qquad\left.  +\left\langle
b_{j}^{\dagger}b_{j+s}^{\dagger}b_{k}^{\dagger}b_{k+s}^{\dagger}G\left(
\mathbf{c}\right)  ^{\dagger N-2}\phi|HG\left(  \mathbf{c}\right)  ^{\dagger
}\phi\right\rangle \right\}
\end{align}
and
\begin{align}
\partial_{c_{j}c_{k}}^{2}S_{N}\left(  \mathbf{c}\right)   &  =N^{2}\left\{
\left\langle b_{j}^{\dagger}b_{j+s}^{\dagger}G\left(  \mathbf{c}\right)
^{\dagger N-1}\phi|b_{k}^{\dagger}b_{k+s}^{\dagger}G\left(  \mathbf{c}\right)
^{\dagger N-1}\right\rangle +\left\langle b_{k}^{\dagger}b_{k+s}^{\dagger
}G\left(  \mathbf{c}\right)  ^{\dagger N-1}\phi|b_{j}^{\dagger}b_{j+s}%
^{\dagger}G\left(  \mathbf{c}\right)  ^{\dagger N-1}\right\rangle \right\}
\nonumber\\
&  +N\left(  N-1\right)  \left\{  \left\langle G\left(  \mathbf{c}\right)
^{\dagger N}\phi|b_{j}^{\dagger}b_{j+s}^{\dagger}b_{k}^{\dagger}%
b_{k+s}^{\dagger}G\left(  \mathbf{c}\right)  ^{\dagger N-2}\phi\right\rangle
\right. \nonumber\\
&  \qquad\qquad\qquad\qquad\qquad\qquad\qquad\qquad\qquad\left.  +\left\langle
b_{j}^{\dagger}b_{j+s}^{\dagger}b_{k}^{\dagger}b_{k+s}^{\dagger}G\left(
\mathbf{c}\right)  ^{\dagger N-2}\phi|G\left(  \mathbf{c}\right)  ^{\dagger
}\phi\right\rangle \right\}  .
\end{align}
Observing
\begin{align}
&  \sum_{1\leq k\leq s}\left\{  \left\langle b_{j}^{\dagger}b_{j+s}^{\dagger
}G\left(  \mathbf{c}\right)  ^{\dagger N-1}\phi|Hb_{k}^{\dagger}%
b_{k+s}^{\dagger}G\left(  \mathbf{c}\right)  ^{\dagger N-1}\right\rangle
+\left\langle b_{k}^{\dagger}b_{k+s}^{\dagger}G\left(  \mathbf{c}\right)
^{\dagger N-1}\phi|Hb_{j}^{\dagger}b_{j+s}^{\dagger}G\left(  \mathbf{c}%
\right)  ^{\dagger N-1}\right\rangle \right\}  c_{k}\nonumber\\
&  =\left\langle b_{j}^{\dagger}b_{j+s}^{\dagger}G\left(  \mathbf{c}\right)
^{\dagger N-1}\phi|HG\left(  \mathbf{c}\right)  ^{\dagger N}\right\rangle
+\left\langle G\left(  \mathbf{c}\right)  ^{\dagger N}\phi|Hb_{j}^{\dagger
}b_{j+s}^{\dagger}G\left(  \mathbf{c}\right)  ^{\dagger N-1}\right\rangle
=\frac{1}{N}\partial_{c_{j}}^{1}\mathcal{H}\left(  \mathbf{c}\right)  ,
\end{align}
and
\begin{align}
&  \sum_{1\leq k\leq s}\left\{  \left\langle G\left(  \mathbf{c}\right)
^{\dagger N}\phi|Hb_{j}^{\dagger}b_{j+s}^{\dagger}b_{k}^{\dagger}%
b_{k+s}^{\dagger}G\left(  \mathbf{c}\right)  ^{\dagger N-2}\phi\right\rangle
+\left\langle b_{j}^{\dagger}b_{j+s}^{\dagger}b_{k}^{\dagger}b_{k+s}^{\dagger
}G\left(  \mathbf{c}\right)  ^{\dagger N-2}\phi|HG\left(  \mathbf{c}\right)
^{\dagger}\phi\right\rangle \right\}  c_{k}\nonumber\\
&  =\left\langle G\left(  \mathbf{c}\right)  ^{\dagger N}\phi|Hb_{j}^{\dagger
}b_{j+s}^{\dagger}G\left(  \mathbf{c}\right)  ^{\dagger N-1}\phi\right\rangle
+\left\langle b_{j}^{\dagger}b_{j+s}^{\dagger}G\left(  \mathbf{c}\right)
^{\dagger N-1}\phi|HG\left(  \mathbf{c}\right)  ^{\dagger N}\phi\right\rangle
=\frac{1}{N}\partial_{c_{j}}^{1}\mathcal{H}\left(  \mathbf{c}\right)
\end{align}
one obtains
\begin{equation}
\sum_{1\leq k\leq s}\partial_{c_{j}c_{k}}^{2}\mathcal{H}\left(  \mathbf{c}%
\right)  c_{k}=\left(  2N-1\right)  \partial_{c_{j}}^{1}\mathcal{H}\left(
\mathbf{c}\right)  .
\end{equation}
The normalization term similarly gives
\begin{align}
&  \sum_{1\leq k\leq s}\left\{  \left\langle b_{j}^{\dagger}b_{j+s}^{\dagger
}G\left(  \mathbf{c}\right)  ^{\dagger N-1}\phi|b_{k}^{\dagger}b_{k+s}%
^{\dagger}G\left(  \mathbf{c}\right)  ^{\dagger N-1}\right\rangle
+\left\langle b_{k}^{\dagger}b_{k+s}^{\dagger}G\left(  \mathbf{c}\right)
^{\dagger N-1}\phi|b_{j}^{\dagger}b_{j+s}^{\dagger}G\left(  \mathbf{c}\right)
^{\dagger N-1}\right\rangle \right\}  c_{k}\nonumber\\
&  =\left\langle b_{j}^{\dagger}b_{j+s}^{\dagger}G\left(  \mathbf{c}\right)
^{\dagger N-1}\phi|G\left(  \mathbf{c}\right)  ^{\dagger N}\right\rangle
+\left\langle G\left(  \mathbf{c}\right)  ^{\dagger N}\phi|b_{j}^{\dagger
}b_{j+s}^{\dagger}G\left(  \mathbf{c}\right)  ^{\dagger N-1}\right\rangle
=\frac{1}{N}\partial_{c_{j}}^{1}S_{N}\left(  \mathbf{c}\right)  ,
\end{align}
and
\begin{align}
&  \sum_{1\leq k\leq s}\left\{  \left\langle G\left(  \mathbf{c}\right)
^{\dagger N}\phi|b_{j}^{\dagger}b_{j+s}^{\dagger}b_{k}^{\dagger}%
b_{k+s}^{\dagger}G\left(  \mathbf{c}\right)  ^{\dagger N-2}\phi\right\rangle
+\left\langle b_{j}^{\dagger}b_{j+s}^{\dagger}b_{k}^{\dagger}b_{k+s}^{\dagger
}G\left(  \mathbf{c}\right)  ^{\dagger N-2}\phi|G\left(  \mathbf{c}\right)
^{\dagger}\phi\right\rangle \right\}  c_{k}\nonumber\\
&  =\left\langle G\left(  \mathbf{c}\right)  ^{\dagger N}\phi|b_{j}^{\dagger
}b_{j+s}^{\dagger}G\left(  \mathbf{c}\right)  ^{\dagger N-1}\phi\right\rangle
+\left\langle b_{j}^{\dagger}b_{j+s}^{\dagger}G\left(  \mathbf{c}\right)
^{\dagger N-1}\phi|G\left(  \mathbf{c}\right)  ^{\dagger N}\phi\right\rangle
=\frac{1}{N}\partial_{c_{j}}^{1}S_{N}\left(  \mathbf{c}\right)  .
\end{align}
producing
\begin{equation}
\sum_{1\leq k\leq s}\partial_{c_{j}c_{k}}^{2}S_{N}\left(  \mathbf{c}\right)
c_{k}=\left(  2N-1\right)  \partial_{c_{j}}^{1}S_{N}\left(  \mathbf{c}\right)
.
\end{equation}
Thus
\begin{equation}
\sum_{1\leq k\leq s}\left\{  \partial_{c_{j}c_{k}}^{2}\mathcal{H}\left(
\mathbf{c}\right)  -\mathcal{E}\left(  \mathbf{c}\right)  \partial_{c_{j}%
c_{k}}^{2}S_{N}\left(  \mathbf{c}\right)  \right\}  c_{k}=\left(  2N-1\right)
\left\{  \partial_{c_{j}}^{1}\mathcal{H}\left(  \mathbf{c}\right)
-\mathcal{E}\left(  \mathbf{c}\right)  \partial_{c_{j}}^{1}S_{N}\left(
\mathbf{c}\right)  \right\}  ,
\end{equation}
which can be written in matrix notation as
\begin{equation}
L\left(  \mathbf{c}\right)  \mathbf{c}\equiv\left[  \partial_{\mathbf{c}}%
^{2}\mathcal{H}\left(  \mathbf{c}\right)  -\mathcal{E}\left(  \mathbf{c}%
\right)  \partial_{\mathbf{c}}^{2}S_{N}\left(  \mathbf{c}\right)  \right]
\mathbf{c=}\left(  2N-1\right)  \partial_{\mathbf{c}}^{1}\mathcal{E}\left(
\mathbf{c}\right)  , \label{gradrel}%
\end{equation}
which clearly shows that
\begin{equation}
\mathbf{c}^{t}L\left(  \mathbf{c}\right)  \mathbf{c=\mathbf{c}}^{t}%
\mathbf{\partial_{\mathbf{c}}^{2}\mathcal{E}\left(  \mathbf{c}\right)
\mathbf{c}}=\mathbf{0.} \label{cd2c}%
\end{equation}


At stationary points, $\mathbf{c}_{\#}$, one can observe from eq. $\left(
\ref{Hess}\right)  $ that
\begin{equation}
\mathbf{\partial_{\mathbf{c}}^{2}\mathcal{E}\left(  \mathbf{c}_{\#}\right)
=}\frac{1}{S_{N}\left(  \mathbf{c}_{\#}\right)  }L\left(  \mathbf{c}%
_{\#}\right)
\end{equation}
and if $\mathbf{c}_{\#}$ is a local minimum then
\begin{equation}
L\left(  \mathbf{c}_{\#}\right)  \geq\mathbf{0.}%
\end{equation}
Hence at a local minimum eq. $\left(  \ref{cd2c}\right)  $ shows that
$\mathbf{c}_{\#}$ is an eigenvector corresponding to the lowest eigenvalue of
$L\left(  \mathbf{c}_{\#}\right)  $ i.e.%
\begin{equation}
L\left(  \mathbf{c}_{\#}\right)  \mathbf{c}_{\#}=0. \label{loweig}%
\end{equation}


If $\mathbf{c}_{\#}$ is the lowest eigenvalue of $L\left(  \mathbf{c}%
_{\#}\right)  $ and $L\left(  \mathbf{c}_{\#}\right)  \geq0$ then as
\begin{equation}
L\left(  \mathbf{c}_{\#}\right)  =S_{N}\left(  \mathbf{c}_{\#}\right)
\partial_{\mathbf{c}}^{2}\mathcal{E}\left(  \mathbf{c}_{\#}\right)
+\partial_{\mathbf{c}}^{1}\mathcal{E}\left(  \mathbf{c}_{\#}\right)
\otimes\partial_{\mathbf{c}}^{1}S_{N}\left(  \mathbf{c}_{\#}\right)
+\partial_{\mathbf{c}}^{1}S_{N}\left(  \mathbf{c}_{\#}\right)  \otimes
\partial_{\mathbf{c}}^{1}\mathcal{E}\left(  \mathbf{c}_{\#}\right)
\end{equation}
one has
\begin{equation}
L\left(  \mathbf{c}_{\#}\right)  \mathbf{c}_{\#}=S_{N}\left(  \mathbf{c}%
_{\#}\right)  \partial_{\mathbf{c}}^{2}\mathcal{E}\left(  \mathbf{c}%
_{\#}\right)  \mathbf{c}_{\#}+2NS_{N}\left(  \mathbf{c}_{\#}\right)
\partial_{\mathbf{c}}^{1}\mathcal{E}\left(  \mathbf{c}_{\#}\right)  =0,
\end{equation}
which implies
\begin{equation}
\frac{1}{2N}\partial_{\mathbf{c}}^{2}\mathcal{E}\left(  \mathbf{c}%
_{\#}\right)  \mathbf{c}_{\#}+\partial_{\mathbf{c}}^{1}\mathcal{E}\left(
\mathbf{c}_{\#}\right)  =0. \label{stat1}%
\end{equation}
If the second order approximation $\mathcal{E}_{2}$
\begin{equation}
\mathcal{E}\left(  \mathbf{c}_{\#}+\mathbf{\delta}\right)  \simeq
\mathcal{E}_{2}\left(  \mathbf{\delta}\right)  =\mathcal{E}\left(
\mathbf{c}_{\#}\right)  +\partial_{\mathbf{c}}^{1}\mathcal{E}\left(
\mathbf{c}_{\#}\right)  ^{t}\mathbf{\delta+}\frac{1}{2}\mathbf{\delta}%
^{t}\partial_{\mathbf{c}}^{2}\mathcal{E}\left(  \mathbf{c}_{\#}\right)
\mathbf{\delta}%
\end{equation}
to $\mathcal{E}$ about the point $\mathbf{c}_{\#}$ is stationary at the point
$\mathbf{\delta}_{\#}$ then
\begin{equation}
\partial_{\mathbf{\delta}}\mathcal{E}_{2}\left(  \mathbf{\delta}_{\#}\right)
=\partial_{\mathbf{c}}^{1}\mathcal{E}\left(  \mathbf{c}_{\#}\right)
^{t}+\mathbf{\delta}^{t}\partial_{\mathbf{c}}^{2}\mathcal{E}\left(
\mathbf{c}_{\#}\right)  =0
\end{equation}
and hence
\begin{equation}
\partial_{\mathbf{c}}^{1}\mathcal{E}\left(  \mathbf{c}_{\#}\right)
+\partial_{\mathbf{c}}^{2}\mathcal{E}\left(  \mathbf{c}_{\#}\right)
\mathbf{\delta}=0 \label{stat2}%
\end{equation}
Thus
\begin{equation}
\mathbf{\delta}_{\#}=\frac{1}{2N}\mathbf{c}_{\#}%
\end{equation}
is a solution to eq. $\left(  \ref{stat2}\right)  $ by observing eq. $\left(
\ref{stat1}\right)  .$ However as $\mathcal{E}$ is a homogenous function of
degree $0$
\begin{equation}
\mathcal{E}\left(  \lambda\mathbf{c}\right)  =\mathcal{E}\left(
\mathbf{c}\right)  \quad\forall\lambda\in\mathbb{R}%
\end{equation}
and
\begin{equation}
\mathcal{E}\left(  \mathbf{c}_{\#}+\mathbf{\delta}_{\#}\right)  =\mathcal{E}%
\left(  \mathbf{c}_{\#}+\frac{1}{2N}\mathbf{c}_{\#}\right)  =\mathcal{E}%
\left(  \mathbf{c}_{\#}\right)  .
\end{equation}
This means $\mathbf{\delta}_{\#}=\mathbf{0}$ also is also a stationary point
of $\mathcal{E}_{2}\left(  \mathbf{\delta}_{\#}\right)  $ and hence eq.
$\left(  \ref{stat2}\right)  $ shows that
\begin{equation}
\partial_{\mathbf{c}}^{1}\mathcal{E}\left(  \mathbf{c}_{\#}\right)
=\mathbf{0,}%
\end{equation}
i.e. $\mathbf{c}_{\#}$ is stationary point of $\mathcal{E}.$ At stationary
points of $\mathcal{E}$
\begin{equation}
\mathbf{\partial_{\mathbf{c}}^{2}\mathcal{E}\left(  \mathbf{c}_{\#}\right)
=}\frac{1}{S_{N}\left(  \mathbf{c}_{\#}\right)  }L\left(  \mathbf{c}%
_{\#}\right)
\end{equation}
and by assumption $L\left(  \mathbf{c}_{\#}\right)  \geq0$ thus one can
conclude that $\mathbf{c}_{\#}$ is a local minimum. We can summarize the
preceding by the theorem

\begin{theorem}
$\mathbf{c}_{\#}$ is a local minimum of $\mathcal{E}$ if and only if
\begin{equation}
L\left(  \mathbf{c}_{\#}\right)  \geq0\qquad\&\qquad L\left(  \mathbf{c}%
_{\#}\right)  \mathbf{c}_{\#}=0.
\end{equation}

\end{theorem}

An iterative minimization scheme can be based on this theorem by considering
the first order approximation
\begin{equation}
\mathcal{E}\left(  \mathbf{c}+\mathbf{\delta}\right)  \simeq\mathcal{E}%
_{1}\left(  \mathbf{\delta}\right)  =\mathcal{E}\left(  \mathbf{c}\right)
+\partial_{\mathbf{c}}^{1}\mathcal{E}\left(  \mathbf{c}\right)  ^{t}%
\mathbf{\delta} \label{o1}%
\end{equation}
and the eigenvalue equation
\begin{equation}
L\left(  \mathbf{c}\right)  \mathbf{x}_{j}=\lambda_{j}\mathbf{x}_{j}\mathbf{,}
\label{Lc}%
\end{equation}
where $\left\{  \left.  \mathbf{x}_{j}\right\vert 1\leq j\leq s\right\}  $ are
orthonormal eigenvectors. By taking the scalar product with $\mathbf{c}$ eq.
$\left(  \ref{Lc}\right)  $ becomes by noting eq. $\left(  \ref{gradrel}%
\right)  $
\begin{equation}
\mathbf{c}^{t}L\left(  \mathbf{c}\right)  \mathbf{x}_{j}=\left(  2N-1\right)
\partial_{\mathbf{c}}^{1}\mathcal{E}\left(  \mathbf{c}\right)  ^{t}%
\mathbf{x}_{j}=\lambda_{j}\mathbf{c}^{t}\mathbf{x}_{j}%
\end{equation}
As $\mathbf{c}$ is always perpendicular to $\partial_{\mathbf{c}}%
^{1}\mathcal{E}\left(  \mathbf{c}\right)  $ we obtain
\begin{equation}
\partial_{\mathbf{c}}^{1}\mathcal{E}\left(  \mathbf{c}\right)  ^{t}\left(
\mathbf{x}_{j}-\mathbf{c}\right)  =\frac{1}{\left(  2N-1\right)  }\lambda
_{j}\mathbf{c}^{t}\mathbf{x}_{j}%
\end{equation}
We can always choose an overall phase for $\mathbf{x}_{j}$ so that
$\mathbf{c}^{t}\mathbf{x}_{j}>0.$ If $\lambda_{L}<0$ we can then set
\begin{equation}
\mathbf{\delta}=\alpha\left(  \mathbf{x}_{L}-\mathbf{c}\right)
\end{equation}
in eq. $\left(  \ref{o1}\right)  $, where $\lambda_{L}$ is the lowest
eigenvalue of $L\left(  \mathbf{c}\right)  $ and find a small enough
$\alpha>0$ that decreases the energy around the point $\mathbf{c.}$ Then we
can compute a new point
\begin{equation}
\mathbf{c}_{new}=\alpha\left(  \mathbf{x}_{L}-\mathbf{c}\right)
\end{equation}
and continue iterating until the lowest eigenvalue is equal to zero.

After the optimal $\mathbf{c}_{\#}$ has been found we can then define an
optimal $\mathbf{n}_{\#}$ by
\begin{equation}
n_{\#j}=c_{\#j}^{2}\qquad1\leq j\leq s.
\end{equation}
As the values $\left\{  c_{\#j}\right\}  $ are optimal
\begin{equation}
\frac{\partial\mathbf{\mathcal{E}\left(  \mathbf{c}_{\#}\right)  }}{\partial
c_{j}}=\frac{\partial\mathbf{\mathcal{E}\left(  \mathbf{n}_{\#}\right)  }%
}{\partial n_{j}}\frac{\partial n_{j}}{\partial c_{j}}=2\frac{\partial
\mathbf{\mathcal{E}\left(  \mathbf{n}_{\#}\right)  }}{\partial n_{j}}%
=0\qquad1\leq j\leq s,
\end{equation}
showing that $\left\{  n_{\#j}\right\}  $ are also optimal and satisfy the
positivity and normalization constraints. The phase information contained in
the signs of $\left\{  c_{\#j}\right\}  ,$ although important in finding a
$\mathbf{c}_{\#}$, is not important for our overall optimization procedure as
it is recovered in the optimization of $\mathbf{\theta}$ and $\mathbf{\psi}$.

\section{Energy Variation Procedure \label{procedure}}

\section{Reduction to the IPS \label{IPS}}

The standard basic model in chemistry is the IPS description, as we have noted
in the preceding this is a special case of the AGP state model which relaxes
to the IPS\ for integer occupation numbers i.e.
\begin{align}
N_{j}  &  =N_{j+s}=1,\quad1\leq j\leq N\nonumber\\
N_{j}  &  =N_{j+s}=0,\quad N+1\leq j\leq s
\end{align}
and%
\begin{align}
n_{j}  &  =n_{j+s}=1,\quad1\leq j\leq N\nonumber\\
n_{j}  &  =n_{j+s}=0,\quad N+1\leq j\leq s.
\end{align}
The IPS\ is invariant to the phases $\mathbf{\theta,}$ which can be set at
$\mathbf{\theta=0}$ and thus the only free variables are the GSO's $\left\{
\mathbf{\psi}\right\}  ,$ leading to the energy expression
\begin{align}
\mathcal{E}\left(  \mathbf{\psi}\right)   &  =\mathcal{H}\left(  \mathbf{\psi
}\right)  =\sum_{1\leq j\leq2s}n_{j}\left\langle \psi_{j}\right\vert
h_{1}\left.  \psi_{j}\right\rangle +\frac{1}{2}\sum_{1\leq j,k\leq2s}%
n_{j}n_{k}\left\{  \left(  \left.  \psi_{j}\psi_{k}\right\vert h_{2}\psi
_{j}\psi_{k}\right)  -\left(  \left.  \psi_{j}\psi_{k}\right\vert h_{2}%
\psi_{k}\psi_{j}\right)  \right\} \nonumber\\
&  =Tr\left\{  \sum_{1\leq j\leq2s}n_{j}\left\vert \psi_{j}\right\rangle
\left\langle \psi_{j}\right\vert h_{1jj}\right. \nonumber\\
&  +\left.  \frac{1}{2}\sum_{1\leq j,k\leq2s}n_{j}n_{k}\left\{  \left\vert
\psi_{j}\right\rangle \otimes\left\vert \psi_{k}\right\rangle \left\langle
\psi_{j}\right\vert \otimes\left\langle \psi_{k}\right\vert -\left\vert
\psi_{j}\right\rangle \otimes\left\vert \psi_{k}\right\rangle \left\langle
\psi_{k}\right\vert \otimes\left\langle \psi_{j}\right\vert \right\}  \left\{
h_{2jkjk}-h_{2jkkj}\right\}  \right\}  .
\end{align}
This can be recast as
\begin{align}
&  \mathcal{H}\left(  \mathbf{\psi}\right)  =Tr\left\{  \sum_{1\leq j,k\leq
2s}n_{j}^{\frac{1}{2}}n_{k}^{\frac{1}{2}}\left\vert \psi_{j}\right\rangle
\left\langle \psi_{k}\right\vert h_{1jk}\right. \nonumber\\
&  \qquad\qquad+\frac{1}{2}\sum_{1\leq j,k,l,m\leq2s}n_{j}^{\frac{1}{2}}%
n_{k}^{\frac{1}{2}}n_{l}^{\frac{1}{2}}n_{m}^{\frac{1}{2}}\left\{  \left\vert
\psi_{j}\right\rangle \otimes\left\vert \psi_{k}\right\rangle -\left\vert
\psi_{k}\right\rangle \otimes\left\vert \psi_{j}\right\rangle \right\}
\nonumber\\
&  \qquad\qquad\qquad\qquad\qquad\qquad\left.  \overset{}{\left\{
\left\langle \psi_{l}\right\vert \otimes\left\langle \psi_{m}\right\vert
-\left\langle \psi_{l}\right\vert \otimes\left\langle \psi_{m}\right\vert
\right\}  \left\{  h_{2jklm}-h_{2jkml}\right\}  }\right\} \nonumber\\
&  \equiv Tr\left\{  \sum_{1\leq j,k\leq2s}\left\vert \psi_{j}\right\rangle
\left\langle \psi_{k}\right\vert K_{1jk}\right. \nonumber\\
&  +\left.  \frac{1}{2}\sum_{1\leq j,k,l,m\leq2s}\left\{  \left\vert \psi
_{j}\right\rangle \otimes\left\vert \psi_{k}\right\rangle \left\langle
\psi_{l}\right\vert \otimes\left\langle \psi_{m}\right\vert -\left\vert
\psi_{j}\right\rangle \otimes\left\vert \psi_{k}\right\rangle \left\langle
\psi_{m}\right\vert \otimes\left\langle \psi_{l}\right\vert \right\}
K_{2jklm}\right\}  ,
\end{align}
where
\begin{equation}
K_{1jk}=n_{j}^{\frac{1}{2}}n_{k}^{\frac{1}{2}}h_{1jk}%
\end{equation}
and
\begin{equation}
K_{2jklm}=n_{j}^{\frac{1}{2}}n_{k}^{\frac{1}{2}}n_{l}^{\frac{1}{2}}%
n_{m}^{\frac{1}{2}}\left\{  h_{2jklm}-h_{2jkml}\right\}  . \label{defk2}%
\end{equation}
One should note that $K_{2jklm}$ is hermitian
\begin{equation}
K_{2jkml}=K_{2mljk}^{\ast}%
\end{equation}
and antisymmetric wrt interchange of the subscripts $j$ $\leftrightarrow$ $k$
and $l$ $\leftrightarrow$ $m$ i.e.
\begin{equation}
K_{2jkml}=-K_{2jklm},K_{2kjlm}=-K_{2kjlm}\quad\text{and }K_{2kjml}=K_{2jklm}.
\end{equation}
This can be seen from definition eq. $\left(  \ref{defk2}\right)  $,
\begin{align}
h_{2abcd}  &  =\int\psi_{a}\left(  \mathbf{x}_{1}\right)  ^{\ast}\psi
_{b}\left(  \mathbf{x}_{2}\right)  ^{\ast}\frac{1}{\left\Vert \mathbf{r}%
_{1}\mathbf{-r}_{2}\right\Vert }\psi_{c}\left(  \mathbf{x}_{1}\right)
\psi_{d}\left(  \mathbf{x}_{2}\right)  d\mathbf{x}_{1}d\mathbf{x}%
_{2}\nonumber\\
&  =\left[  \int\psi_{b}\left(  \mathbf{x}_{1}\right)  \psi_{a}\left(
\mathbf{x}_{2}\right)  \frac{1}{\left\Vert \mathbf{r}_{1}\mathbf{-r}%
_{2}\right\Vert }\psi_{d}\left(  \mathbf{x}_{1}\right)  ^{\ast}\psi_{c}\left(
\mathbf{x}_{2}\right)  ^{\ast}d\mathbf{x}_{1}d\mathbf{x}_{2}\right]  ^{\ast
}=h_{2dcba}^{\ast}%
\end{align}
and
\begin{align}
h_{2abcd}  &  =\int\psi_{a}\left(  \mathbf{x}_{1}\right)  ^{\ast}\psi
_{b}\left(  \mathbf{x}_{2}\right)  ^{\ast}\frac{1}{\left\Vert \mathbf{r}%
_{1}\mathbf{-r}_{2}\right\Vert }\psi_{c}\left(  \mathbf{x}_{1}\right)
\psi_{d}\left(  \mathbf{x}_{2}\right)  d\mathbf{x}_{1}d\mathbf{x}%
_{2}\nonumber\\
&  =\int\psi_{b}\left(  \mathbf{x}_{1}\right)  ^{\ast}\psi_{a}\left(
\mathbf{x}_{2}\right)  ^{\ast}\frac{1}{\left\Vert \mathbf{r}_{1}%
\mathbf{-r}_{2}\right\Vert }\psi_{d}\left(  \mathbf{x}_{1}\right)  \psi
_{c}\left(  \mathbf{x}_{2}\right)  d\mathbf{x}_{1}d\mathbf{x}_{2}=h_{2badc}.
\end{align}


The functional derivatives wrt $\left\{  \psi_{\mu}^{\ast};1\leq\mu
\leq2s\right\}  $ are given by%
\begin{align}
\frac{\delta\mathcal{H}\left(  \mathbf{\psi,n}\right)  }{\delta\psi_{\mu
}^{\ast}}  &  =\sum_{1\leq j\leq2s}\left\vert \psi_{j}\right\rangle K_{1j\mu
}+Tr\left\{  \frac{1}{2}\sum_{1\leq j,k,l\leq2s}\left\{  \left\vert \psi
_{j}\right\rangle \left\{  \left\vert \psi_{k}\right\rangle \left\langle
\psi_{l}\right\vert \right\}  K_{2kjjl\mu}+\left\vert \psi_{j}\right\rangle
\left\vert \psi_{k}\right\rangle \left\langle \psi_{l}\right\vert K_{2jk\mu
l}\right\}  \right\} \nonumber\\
&  =\sum_{1\leq j\leq2s}\left\vert \psi_{j}\right\rangle K_{1j\mu}+Tr\left\{
\sum_{1\leq j,k,l\leq2s}\left\vert \psi_{j}\right\rangle \left\{  \left\vert
\psi_{k}\right\rangle \left\langle \psi_{l}\right\vert \right\}  K_{2jk\mu
l}\right\}  \equiv\sum_{1\leq j\leq2s}\left\vert \psi_{j}\right\rangle
F\left(  \mathbf{\psi}\right)  _{j\mu},
\end{align}
where
\begin{equation}
F\left(  \mathbf{\psi}\right)  _{j\mu}=K_{1j\mu}+Tr\left\{  \sum_{1\leq
k,l\leq2s}\left\vert \psi_{k}\right\rangle \left\langle \psi_{l}\right\vert
K_{2jk\mu l}\right\}  =K_{1j\mu}+\sum_{1\leq k\leq2s}K_{2jk\mu k}.
\end{equation}
This defines a Fock operator by
\begin{equation}
F\left(  \mathbf{\psi}\right)  \left(  \left\vert \psi_{1}\right\rangle
,\cdots,\left\vert \psi_{2s}\right\rangle \right)  =\left(  \left\vert
\psi_{1}\right\rangle ,\cdots,\left\vert \psi_{2s}\right\rangle \right)
\mathbb{F}\left(  \mathbf{\psi}\right)  ,
\end{equation}
where $\mathbb{F}\left(  \mathbf{\psi}\right)  $ has matrix elements $\left\{
F\left(  \mathbf{\psi}\right)  _{j\mu}\right\}  $ wrt the basis $\left\{
\left\vert \psi_{j}\right\rangle ;1\leq j\leq2s\right\}  .$

The functional derivatives wrt $\left\{  \psi_{\mu};1\leq\mu\leq2s\right\}  $
are given by%
\begin{equation}
r\left\{  \sum_{1\leq j,k\leq2s}\left\vert \psi_{j}\right\rangle \left\langle
\psi_{k}\right\vert K_{1jk}+\frac{1}{2}\sum_{1\leq j,k,l,m\leq2s}\left[
\left\vert \psi_{j}\right\rangle \otimes\left\vert \psi_{k}\right\rangle
\right]  \left[  \left\langle \psi_{l}\right\vert \otimes\left\langle \psi
_{m}\right\vert \right]  K_{2jklm}\right\}
\end{equation}
\begin{align}
\frac{\delta\mathcal{H}\left(  \mathbf{\psi,n}\right)  }{\delta\psi_{\mu}}  &
=\sum_{1\leq j\leq2s}K_{1\mu j}\left\langle \psi_{\mu}\right\vert +Tr\left\{
\frac{1}{2}\sum_{1\leq j,l,m\leq2s}\left\vert \psi_{j}\right\rangle \left[
\left\langle \psi_{l}\right\vert \otimes\left\langle \psi_{m}\right\vert
\right]  K_{2\mu jlm}\right\} \nonumber\\
&  \sum_{1\leq j,l,m\leq2s}\left[  \left\langle \psi_{l}\right\vert
\otimes\left[  \left\vert \psi_{j}\right\rangle \right]  \left\langle \psi
_{m}\right\vert \right]  K_{2j\mu lm}\\
&  =\left\{  K_{1\mu\mu}+Tr\left\{  \sum_{1\leq k\leq2s}K_{2\mu k\mu
k}\right\}  \right\}  \left\langle \psi_{\mu}\right\vert \equiv\left\langle
\psi_{\mu}\right\vert F\left(  \mathbf{\psi}\right)  _{\mu\mu}.
\end{align}


\section{Discussion\label{Discussion}}

We have extended the IPS one particle model of molecular quantum mechanics
described by single determinantal wavefunctions to a more accurate correlated
one particle model described by AGP\ wavefunctions for an even number of
electrons and GAGP wavefunctions for an odd number. All of these wavefunctions
are determined by a set of one particle occupation numbers and a set of GSO's.
In the determinantal case the occupation numbers are one or zero, while in the
AGP case they must be at least doubly degenerate and for the GAGP case there
must be at most one non degenerate integer occupation number and the rest at
least doubly degenerate.

The occupation numbers and GSO's totally determine the many electron state, a
necessary criterion for a one particle model. In the AGP\ case the GSO's
produce weighted Probability Distributions (PD's)\ that produce pictorial
models of a molecules that generalize the standard ones used in chemistry in
that they contain the effects of correlation. It is true that multi
configurational wavefunctions also produce weighted Probability Distributions
(PD's) but they are not determined by them in contrast to the AGP\ case.

Another important aspect of the IPS model is the role of energies of the one
particle orbitals, which are given by the eigenvalues of the Fock operator.
When these are combined with the shape of PD's one obtains a description of
chemical properties i.e. reactivity and reaction mechanisms. Physically these
energies, are the poles of a first order approximation to the electron
propagator and correspond to the ionization potentials of the electrons from
the occupied orbitals. Mathematically they reflect the effect of the
orthonormality constraints in the IP-SCF\ procedure i.e. they give the rate of
change of the total energy with respect to changes in the normalization of the
corresponding orbitals. The aufbau principle is also based on these
eigenvalues i.e. the optimized IPS\ is in general composed of the orbitals
corresponding to most negative ionization potentials.

These properties are also generalized by the self consistent AGP-Fock
operator. However there some important differences, the IPS is invariant to
phase changes of the eigenvectors of the SCF-Fock operator, while the AGP-Fock
operator and the total energy are \emph{not }invariant to such changes. Hence
one must perform an independent optimization of these phases in order to
determine the best AGP approximation to the true many electron ground state.
This phase dependence is a reflection of the fact the AGP FORDO does not
determine a unique AGP\ state, but only an equivalence class of states. By
optimizing over the phase for fixed reference FORDO one obtains an explicit
Density Matrix Function (DMF)\ that can be used as a model functional for
DMFT. We have also shown, by analyzing the total energy in terms of the
eigenvalues of the AGP-Fock operator, that a generalized aufbau principle
holds i.e. the lowest AGP-Fock operator eigenvalues are associated with the
most highly occupied orbitals.

The AGP-Fock operator eigenvalues again reflect the effect of the
orthonormality constraints, which we show are sufficient to enforce
orthogonality between the orbitals in \emph{both }the IP and AGP state cases.
This unconventional approach to the derivation of the IP-SCF equations is a
necessity in the SCF AGP\ equations as the standard arguments used in the
IPS\ case do not in this case lead to an eigenvalue equation.

The final variables that need to be optimized in the AGP\ case are the
occupation numbers of the FORDO, which we can express in terms of the
occupation numbers of the geminal using the monotonic $1-1$ correspondence
shown by \cite{Coleman} and subsequent angular functions in order to maintain
the normalization and positivity constraints. The relaxation of the occupation
numbers from integer values and the determination of optimal phases reflect
the correlation effects.

We have not examined the relationship of the AGP-Fock operator to the poles of
the one electron propagator in any rigorous fashion, nor have we considered
the explicit form of the derivatives of the total energy $\left\{
\frac{\partial\mathcal{E}}{\partial N_{j}};1\leq j\leq s\right\}  $ with
respect to the occupation numbers of the FORDO in order to study Janaks
Theorem \cite{Janak}.

\section{Acknowledgements}

\appendix

\section{One Electron Basis\label{basis}}

\section{Example of Linear Fit\label{LinfitEX}}

consider%
\begin{equation}
\mathbf{N}_{1}=\left(
\begin{array}
[c]{c}%
1\\
.9\\
.8\\
.1\\
.2\\
0
\end{array}
\right)  =\left(
\begin{array}
[c]{c}%
1\\
1\\
1\\
0\\
0\\
0
\end{array}
\right)  +\left(
\begin{array}
[c]{c}%
0\\
-.1\\
-.2\\
.1\\
.2\\
0
\end{array}
\right)  =\mathbf{N}_{0}+\delta\mathbf{N}_{1}%
\end{equation}%
\begin{equation}
\mathbf{n}_{1}=\mathcal{N}^{-1}\left(
\begin{array}
[c]{c}%
1\\
.9\\
.8\\
0\\
.2\\
.1
\end{array}
\right)  =\left(
\begin{array}
[c]{c}%
\frac{1}{3}\\
\frac{1}{3}\\
\frac{1}{3}\\
0\\
0\\
0
\end{array}
\right)  +\left[  d\mathcal{N}\left(
\begin{array}
[c]{c}%
\frac{1}{3}\\
\frac{1}{3}\\
\frac{1}{3}\\
0\\
0\\
0
\end{array}
\right)  \right]  ^{-1}\left(
\begin{array}
[c]{c}%
0\\
-.1\\
-.2\\
.1\\
.2\\
0
\end{array}
\right)
\end{equation}%
\[
\mathbf{n}_{2}=\mathcal{N}^{-1}\left(
\begin{array}
[c]{c}%
.9\\
.8\\
.8\\
.1\\
.2\\
.2
\end{array}
\right)  =\left(
\begin{array}
[c]{c}%
\frac{1}{3}\\
\frac{1}{3}\\
\frac{1}{3}\\
0\\
0\\
0
\end{array}
\right)  +\left[  d\mathcal{N}\left(
\begin{array}
[c]{c}%
\frac{1}{3}\\
\frac{1}{3}\\
\frac{1}{3}\\
0\\
0\\
0
\end{array}
\right)  \right]  ^{-1}\left(
\begin{array}
[c]{c}%
0\\
-.1\\
-.2\\
.1\\
.2\\
0
\end{array}
\right)
\]%
\begin{equation}
+\left[  d\mathcal{N}\left(  \left(
\begin{array}
[c]{c}%
\frac{1}{3}\\
\frac{1}{3}\\
\frac{1}{3}\\
0\\
0\\
0
\end{array}
\right)  +\left[  d\mathcal{N}\left(
\begin{array}
[c]{c}%
\frac{1}{3}\\
\frac{1}{3}\\
\frac{1}{3}\\
0\\
0\\
0
\end{array}
\right)  \right]  ^{-1}\left(
\begin{array}
[c]{c}%
0\\
-.1\\
-.2\\
.1\\
.2\\
0
\end{array}
\right)  \right)  \right]  ^{-1}\left(
\begin{array}
[c]{c}%
-.1\\
-.1\\
0\\
.1\\
0\\
.1
\end{array}
\right)
\end{equation}%
\begin{equation}
\equiv\mathbf{n}_{0}+\left[  d\mathcal{N}\left(  \mathbf{n}_{0}\right)
\right]  ^{-1}\left(  \delta\mathbf{N}_{1}\right)  +\left[  d\mathcal{N}%
\left(  \mathbf{n}_{0}+\left[  d\mathcal{N}\left(  \mathbf{n}_{0}\right)
\right]  ^{-1}\left(  \delta\mathbf{N}_{1}\right)  \right)  \right]
^{-1}\left(  \left(  \delta\mathbf{N}_{2}\right)  \right)
\end{equation}


\section{Variational Equations\label{variational}}

\subsection{Derivatives}

The $2N$-AGP\ energy functional in terms of the CGSO's $\left\{  \left.
\psi_{j}\right\vert 1\leq j\leq2s\right\}  $, the geminal occupation numbers
$\left\{  \left.  n_{j}\right\vert 1\leq j\leq s\right\}  $ and the phases
$\left\{  \left.  \theta_{j}\right\vert 1\leq j\leq s\right\}  $ can be
expressed as the quotient%
\begin{align}
&  \mathcal{E}\left(  \mathbf{n},\mathbf{\psi},\mathbf{\theta}\right)
=\frac{\mathcal{H}\left(  \mathbf{n},\mathbf{\psi},\mathbf{\theta}\right)
}{S_{N}\left(  \mathbf{n}\right)  }\nonumber\\
&  =\frac{1}{S_{N}\left(  \mathbf{n}\right)  }\left\{  \sum_{1\leq j\leq
s}n_{j}\left\{  \left\langle \psi_{j}\left\vert h_{1}\psi_{j}\right.
\right\rangle +\left\langle \psi_{j+s}\left\vert h_{1}\psi_{j+s}\right.
\right\rangle +\psi_{j}\psi_{j+s}|\left\vert \psi_{j}\psi_{j+s}\right\rangle
\right\}  \right. \nonumber\\
&  +\sum_{1\leq j,k\leq s}n_{j}^{\frac{1}{2}}n_{k}^{\frac{1}{2}}e^{i\left(
\theta_{j}-\theta_{k}\right)  }S_{N-1}\left(  \jmath k\right)  \psi_{j}%
\psi_{j+s}|\left\vert \psi_{k}\psi_{k+s}\right\rangle \left.  +\frac{1}{2}%
\sum_{1\leq j,k\leq s}n_{j}n_{k}S_{N-2}\left(  \jmath k\right)  \left\langle
\psi_{j}\psi_{k}|||\psi_{j}\psi_{k}\right\rangle \right\}  .
\end{align}
The first order variation $\delta^{1}\mathcal{E}\left(  \mathbf{n}%
,\mathbf{\psi},\mathbf{\theta}\right)  $ of the energy is
\begin{equation}
\delta^{1}\mathcal{E}\left(  \mathbf{n},\mathbf{\psi},\mathbf{\theta}\right)
=\left(
\begin{array}
[c]{cccc}%
\partial_{\mathbf{n}}^{1}\mathcal{E}\left(  \mathbf{n},\mathbf{\psi
},\mathbf{\theta}\right)  & \partial_{\mathbf{\psi}^{\ast}}^{1}\mathcal{E}%
\left(  \mathbf{n},\mathbf{\psi},\mathbf{\theta}\right)  & \partial
_{\mathbf{\psi}}^{1}\mathcal{E}\left(  \mathbf{n},\mathbf{\psi},\mathbf{\theta
}\right)  & \partial_{\mathbf{\theta}}^{1}\mathcal{E}\left(  \mathbf{n}%
,\mathbf{\psi},\mathbf{\theta}\right)
\end{array}
\right)  \left(
\begin{array}
[c]{c}%
\delta\mathbf{n}\\
\delta\mathbf{\psi}^{\ast}\\
\delta\mathbf{\psi}\\
\delta\mathbf{\theta}%
\end{array}
\right)  ,
\end{equation}
where
\begin{subequations}
\begin{align}
\partial_{\mathbf{n}}^{1}\mathcal{E}\left(  \mathbf{n},\mathbf{\psi
},\mathbf{\theta}\right)   &  =\frac{1}{S_{N}\left(  \mathbf{n}\right)
}\left\{  \partial_{\mathbf{n}}^{1}\mathcal{H}\left(  \mathbf{n},\mathbf{\psi
},\mathbf{\theta}\right)  -\mathcal{E}\left(  \mathbf{n},\mathbf{\psi
},\mathbf{\theta}\right)  \partial_{\mathbf{n}}^{1}S_{N}\left(  \mathbf{n}%
\right)  \right\} \\
\partial_{\mathbf{\psi}^{\ast}}^{1}\mathcal{E}\left(  \mathbf{n},\mathbf{\psi
},\mathbf{\theta}\right)   &  =\frac{1}{S_{N}\left(  \mathbf{n}\right)
}\left\{  \partial_{\mathbf{\psi}^{\ast}}^{1}\mathcal{H}\left(  \mathbf{n}%
,\mathbf{\psi},\mathbf{\theta}\right)  -\mathcal{E}\left(  \mathbf{n}%
,\mathbf{\psi},\mathbf{\theta}\right)  \partial_{\mathbf{\psi}^{\ast}}%
^{1}S_{N}\left(  \mathbf{n}\right)  \right\}  =\left[  \partial_{\mathbf{\psi
}}^{1}\mathcal{E}\left(  \mathbf{n},\mathbf{\psi},\mathbf{\theta}\right)
\right]  ^{\ast}\\
\partial_{\mathbf{\theta}}^{1}\mathcal{E}\left(  \mathbf{n},\mathbf{\psi
},\mathbf{\theta}\right)   &  =\frac{1}{S_{N}\left(  \mathbf{n}\right)
}\left\{  \partial_{\mathbf{\theta}}^{1}\mathcal{H}\left(  \mathbf{n}%
,\mathbf{\psi},\mathbf{\theta}\right)  -\mathcal{E}\left(  \mathbf{n}%
,\mathbf{\psi},\mathbf{\theta}\right)  \partial_{\mathbf{\theta}}^{1}%
S_{N}\left(  \mathbf{n}\right)  \right\}
\end{align}
and
\end{subequations}
\begin{subequations}
\begin{align}
\partial_{\mathbf{\psi}^{\ast}}^{1}\mathcal{E}\left(  \mathbf{n},\mathbf{\psi
},\mathbf{\theta}\right)  \cdot\delta\mathbf{\psi}^{\ast}  &  =\sum_{1\leq
j\leq2s}\partial_{\psi_{j}^{\ast}}^{1}\mathcal{E}\left(  \mathbf{n}%
,\mathbf{\psi},\mathbf{\theta}\right)  \left(  \delta\psi_{j}^{\ast}\right)
=\sum_{1\leq j\leq2s}\int\frac{\delta\mathcal{E}\left(  \mathbf{n}%
,\mathbf{\psi},\mathbf{\theta}\right)  }{\delta\psi_{j}^{\ast}\left(
\mathbf{r,}\zeta\right)  }\delta\psi_{j}^{\ast}\left(  \mathbf{r,}%
\zeta\right)  d\mathbf{r}d\zeta\label{funcder}\\
\partial_{\mathbf{n}}^{1}\mathcal{E}\left(  \mathbf{n},\mathbf{\psi
},\mathbf{\theta}\right)  \cdot\delta\mathbf{n}  &  =\sum_{1\leq j\leq s}%
\frac{\partial\mathcal{E}\left(  \mathbf{n},\mathbf{\psi},\mathbf{\theta
}\right)  }{\partial n_{j}}\delta n_{j}\\
\partial_{\mathbf{\theta}}^{1}\mathcal{E}\left(  \mathbf{n},\mathbf{\psi
},\mathbf{\theta}\right)  \cdot\delta\mathbf{\theta}  &  =\sum_{1\leq j\leq
s}\frac{\partial\mathcal{E}\left(  \mathbf{n},\mathbf{\psi},\mathbf{\theta
}\right)  }{\partial\theta_{j}}\delta\theta_{j}.
\end{align}
If the variations $\left\{  \delta\mathbf{\psi}^{\ast},\delta\mathbf{\psi
}\right\}  $ maintain the orthonormality of the CGSO's i.e.
\end{subequations}
\begin{equation}
\left\langle \left.  \psi_{j}\right\vert \psi_{k}\right\rangle =\delta_{jk}%
\end{equation}
then
\begin{align}
\partial_{\mathbf{\psi}^{\ast}}^{1}\mathcal{E}\left(  \mathbf{n},\mathbf{\psi
},\mathbf{\theta}\right)   &  =\frac{1}{S_{N}\left(  \mathbf{n}\right)
}\partial_{\mathbf{\psi}^{\ast}}^{1}\mathcal{H}\left(  \mathbf{n}%
,\mathbf{\psi},\mathbf{\theta}\right) \\
\partial_{\mathbf{\theta}}^{1}\mathcal{E}\left(  \mathbf{n},\mathbf{\psi
},\mathbf{\theta}\right)   &  =\frac{1}{S_{N}\left(  \mathbf{n}\right)
}\partial_{\mathbf{\theta}}^{1}\mathcal{H}\left(  \mathbf{n},\mathbf{\psi
},\mathbf{\theta}\right)
\end{align}
as $\partial_{\mathbf{\psi}^{\ast}}^{1}S_{N}\left(  \mathbf{n}\right)
=\partial_{\mathbf{\psi}}^{1}S_{N}\left(  \mathbf{n}\right)  =\partial
_{\mathbf{\theta}}^{1}S_{N}\left(  \mathbf{n}\right)  =\mathbf{0}$.

\subsection{Derivation of the Variational Equations for the Geminal Occupation
Numbers}

\subsection{Derivation of the Fock Operator}

In order to obtain derivative of terms in the energy expression it helpful to
utilize matrix elements wrt non antisymmetric states
\begin{equation}
\left\vert \psi_{j}\psi_{k}\right)  =\left\vert \psi_{j}\right)
\otimes\left\vert \psi_{k}\right)  ,
\end{equation}
whose amplitudes are
\begin{equation}
\left(  \left.  \mathbf{r}_{1}\mathbf{,\xi}_{1}\mathbf{r}_{2}\mathbf{,\xi}%
_{2}\right\vert \psi_{j}\psi_{k}\right)  =\psi_{j}\left(  \mathbf{r}%
_{1}\mathbf{,\xi}_{1}\right)  \psi_{k}\left(  \mathbf{r}_{2}\mathbf{,\xi}%
_{2}\right)  ,
\end{equation}
where
\begin{equation}
\psi_{j}\left(  \mathbf{r}_{1}\mathbf{,\xi}_{1}\right)  =\left(  \left.
\mathbf{r}_{1}\mathbf{,\xi}_{1}\right\vert \psi_{j}\right)  =\left\langle
\left.  \mathbf{r}_{1}\mathbf{,\xi}_{1}\right\vert \psi_{j}\right\rangle .
\end{equation}
In terms of these states the unnormalized $2N$-electron AGP\ energy becomes
\begin{align}
&  \mathcal{H}\left(  \mathbf{n},\mathbf{\psi},\mathbf{\theta}\right)
=\mathcal{H}_{1}\left(  \mathbf{n},\mathbf{\psi},\mathbf{\theta}\right)
+\mathcal{H}_{2}\left(  \mathbf{n},\mathbf{\psi},\mathbf{\theta}\right)
=\nonumber\\
&  \sum_{1\leq j\leq s}n_{j}S_{N-1}\left(  \jmath\right)  \left\{  \left(
\left.  \psi_{j}\right\vert h_{1}\psi_{j}\right)  +\left(  \left.  \psi
_{j+s}\right\vert h_{1}\psi_{j+s}\right)  +\left(  \left.  \psi_{j}\psi
_{j+s}\right\vert h_{2}\psi_{j}\psi_{j+s}\right)  -\left(  \left.  \psi
_{j}\psi_{j+s}\right\vert h_{2}\psi_{j+s}\psi_{j}\right)  \right\} \nonumber\\
&  \qquad+\sum_{1\leq j,k\leq s}\left\{  n_{j}^{\frac{1}{2}}n_{k}^{\frac{1}%
{2}}e^{i\left(  \theta_{j}-\theta_{k}\right)  }S_{N-1}\left(  \jmath k\right)
\left\{  \left(  \left.  \psi_{j}\psi_{j+s}\right\vert h_{2}\psi_{k}\psi
_{k+s}\right)  -\left(  \left.  \psi_{j}\psi_{j+s}\right\vert h_{2}\psi
_{k+s}\psi_{k}\right)  \right\}  \right. \nonumber\\
&  \qquad\qquad+\frac{1}{2}\sum_{1\leq j,k\leq2s}n_{j}n_{k}S_{N-2}\left(
\jmath k\right)  \overset{}{\left\{  \left(  \left.  \psi_{j}\psi
_{k}\right\vert h_{2}\psi_{j}\psi_{k}\right)  -\left(  \left.  \psi_{j}%
\psi_{k}\right\vert h_{2}\psi_{k}\psi_{j}\right)  \right\}  }, \label{unnormE}%
\end{align}
where
\begin{equation}
\mathcal{H}_{1}\left(  \mathbf{n},\mathbf{\psi},\mathbf{\theta}\right)
=\sum_{1\leq j\leq2s}n_{j}S_{N-1}\left(  \jmath\right)  \left\langle \left.
\psi_{j}\right\vert h_{1}\psi_{j}\right\rangle
\end{equation}
and
\begin{align}
&  \mathcal{H}_{2}\left(  \mathbf{n},\mathbf{\psi},\mathbf{\theta}\right)
=\sum_{1\leq j\leq2s}n_{j}S_{N-1}\left(  \jmath\right)  \left\{  \left(
\left.  \psi_{j}\psi_{j+s}\right\vert h_{2}\psi_{j}\psi_{j+s}\right)  -\left(
\left.  \psi_{j}\psi_{j+s}\right\vert h_{2}\psi_{j+s}\psi_{j}\right)  \right\}
\nonumber\\
&  \qquad+\sum_{1\leq j,k\leq s}n_{j}^{\frac{1}{2}}n_{k}^{\frac{1}{2}%
}e^{i\left(  \theta_{j}-\theta_{k}\right)  }S_{N-1}\left(  \jmath k\right)
\left\{  \left(  \left.  \psi_{j}\psi_{j+s}\right\vert h_{2}\psi_{k}\psi
_{k+s}\right)  -\left(  \left.  \psi_{j}\psi_{j+s}\right\vert h_{2}\psi
_{k+s}\psi_{k}\right)  \right\} \nonumber\\
&  \qquad\qquad\qquad+\frac{1}{2}\sum_{1\leq j,k\leq2s}n_{j}n_{k}%
S_{N-2}\left(  \jmath k\right)  \overset{}{\left\{  \left(  \left.  \psi
_{j}\psi_{k}\right\vert h_{2}\psi_{j}\psi_{k}\right)  -\left(  \left.
\psi_{j}\psi_{k}\right\vert h_{2}\psi_{k}\psi_{j}\right)  \right\}  }.
\end{align}
The form of $h_{1}$ and $h_{2}$ can be seen from eqs. $\left(  \ref{h1}%
\right)  $ and $\left(  \ref{h2}\right)  .$ The two particle integrals
$\left\{  h_{2abcd}\right\}  $ have the properties
\begin{align}
h_{2abcd}  &  =\int\psi_{a}\left(  \mathbf{x}_{1}\right)  ^{\ast}\psi
_{b}\left(  \mathbf{x}_{2}\right)  ^{\ast}\frac{1}{\left\Vert \mathbf{r}%
_{1}\mathbf{-r}_{2}\right\Vert }\psi_{c}\left(  \mathbf{x}_{1}\right)
\psi_{d}\left(  \mathbf{x}_{2}\right)  d\mathbf{x}_{1}d\mathbf{x}%
_{2}\nonumber\\
&  =\left[  \int\psi_{b}\left(  \mathbf{x}_{1}\right)  \psi_{a}\left(
\mathbf{x}_{2}\right)  \frac{1}{\left\Vert \mathbf{r}_{1}\mathbf{-r}%
_{2}\right\Vert }\psi_{d}\left(  \mathbf{x}_{1}\right)  ^{\ast}\psi_{c}\left(
\mathbf{x}_{2}\right)  ^{\ast}d\mathbf{x}_{1}d\mathbf{x}_{2}\right]  ^{\ast
}=h_{2cdab}^{\ast}%
\end{align}
and
\begin{align}
h_{2abcd}  &  =\int\psi_{a}\left(  \mathbf{x}_{1}\right)  ^{\ast}\psi
_{b}\left(  \mathbf{x}_{2}\right)  ^{\ast}\frac{1}{\left\Vert \mathbf{r}%
_{1}\mathbf{-r}_{2}\right\Vert }\psi_{c}\left(  \mathbf{x}_{1}\right)
\psi_{d}\left(  \mathbf{x}_{2}\right)  d\mathbf{x}_{1}d\mathbf{x}%
_{2}\nonumber\\
&  =\int\psi_{b}\left(  \mathbf{x}_{1}\right)  ^{\ast}\psi_{a}\left(
\mathbf{x}_{2}\right)  ^{\ast}\frac{1}{\left\Vert \mathbf{r}_{1}%
\mathbf{-r}_{2}\right\Vert }\psi_{d}\left(  \mathbf{x}_{1}\right)  \psi
_{c}\left(  \mathbf{x}_{2}\right)  d\mathbf{x}_{1}d\mathbf{x}_{2}=h_{2badc}.
\end{align}
The first order change in the energy due to changes $\left\{  \delta
\mathbf{\psi}^{\ast},\delta\mathbf{\psi}\right\}  $in the GSO's is given by
the Taylors series as
\begin{align}
\delta^{1}\mathcal{H}\left(  \mathbf{n},\mathbf{\psi},\mathbf{\theta}\right)
&  =\left(  \partial_{\mathbf{\psi}^{\ast}}^{1}\mathcal{H}\left(
\mathbf{n},\mathbf{\psi},\mathbf{\theta}\right)  ,\partial_{\mathbf{\psi}}%
^{1}\mathcal{H}\left(  \mathbf{n},\mathbf{\psi},\mathbf{\theta}\right)
\right)  \left(
\begin{array}
[c]{c}%
\delta\mathbf{\psi}^{\ast}\\
\delta\mathbf{\psi}%
\end{array}
\right) \nonumber\\
&  =\left(  \partial_{\mathbf{\psi}^{\ast}}^{1}\mathcal{H}\left(
\mathbf{n},\mathbf{\psi},\mathbf{\theta}\right)  ,\partial_{\mathbf{\psi
}^{\ast}}^{1}\mathcal{H}\left(  \mathbf{n},\mathbf{\psi},\mathbf{\theta
}\right)  ^{\ast}\right)  \left(
\begin{array}
[c]{c}%
\delta\mathbf{\psi}^{\ast}\\
\delta\mathbf{\psi}%
\end{array}
\right) \nonumber\\
&  =2\operatorname{Re}\left\{  \partial_{\mathbf{\psi}^{\ast}}^{1}%
\mathcal{H}\left(  \mathbf{n},\mathbf{\psi},\mathbf{\theta}\right)  ^{\dagger
}\delta\mathbf{\psi}^{\ast}\right\}  ,
\end{align}
where
\begin{equation}
\partial_{\psi_{\mu}^{\ast}}^{1}\mathcal{H}\left(  \mathbf{n},\mathbf{\psi
},\mathbf{\theta}\right)  =\frac{\delta\mathcal{H}\left(  \mathbf{n}%
,\mathbf{\psi},\mathbf{\theta}\right)  }{\delta\psi_{\mu}^{\ast}}.
\end{equation}
It is clear that
\begin{equation}
\partial_{\psi_{\mu}^{\ast}}^{1}\mathcal{H}\left(  \mathbf{n},\mathbf{\psi
},\mathbf{\theta}\right)  =\left[  \partial_{\psi_{\mu}}^{1}\mathcal{H}\left(
\mathbf{n},\mathbf{\psi},\mathbf{\theta}\right)  \right]  ^{\ast}%
\end{equation}
as $\delta^{1}\mathcal{H}\left(  \mathbf{n},\mathbf{\psi},\mathbf{\theta
}\right)  $ must be real valued for all $\left(  \delta\mathbf{\psi}^{\ast
},\delta\mathbf{\psi}\right)  $. In order to find the explicit expressions for
the derivatives of eq. $\left(  \ref{unnormE}\right)  $ one observes that the
one electron functional derivatives are given by
\begin{equation}
\frac{\delta}{\delta\psi_{\mu}^{\ast}}\left\langle \left.  \psi_{a}\right\vert
h_{1}\psi_{b}\right\rangle =h_{1}\left\vert \psi_{b}\right\rangle \delta
_{a\mu}=\left\{  h_{1}\left\vert \psi_{b}\right\rangle \left\langle \psi
_{a}\right\vert \right\}  \left\vert \psi_{\mu}\right\rangle .
\label{1partfuncder_o}%
\end{equation}
After noting that
\begin{equation}
h_{1}\left\vert \psi_{b}\right\rangle \left\langle \psi_{a}\right\vert
=\sum_{1\leq j,k\leq2s}h_{1jk}\left\vert \psi_{j}\right\rangle \left\langle
\left.  \psi_{k}\right\vert \psi_{b}\right\rangle \left\langle \psi
_{a}\right\vert =\sum_{1\leq j\leq2s}h_{1jb}\left\vert \psi_{j}\right\rangle
\left\langle \psi_{a}\right\vert
\end{equation}
eq. $\left(  \ref{1partfuncder_o}\right)  $ can be expanded as
\begin{equation}
\frac{\delta}{\delta\psi_{\mu}^{\ast}}\left\langle \left.  \psi_{a}\right\vert
h_{1}\psi_{b}\right\rangle =\left\{  \sum_{1\leq j,k\leq2s}h_{1jb}\delta
_{ak}\left\vert \psi_{j}\right\rangle \left\langle \psi_{k}\right\vert
\right\}  \left\vert \psi_{\mu}\right\rangle . \label{1partfuncder_ex}%
\end{equation}
Thus
\begin{align}
\frac{\delta}{\delta\psi_{\mu}^{\ast}}\mathcal{H}_{1}\left(  \mathbf{n}%
,\mathbf{\psi},\mathbf{\theta}\right)   &  =\sum_{1\leq a\leq2s}n_{a}%
S_{N-1}\left(  a\right)  \frac{\delta}{\delta\psi_{\mu}^{\ast}}\left\langle
\left.  \psi_{a}\right\vert h_{1}\psi_{a}\right\rangle \\
&  =\left\{  \sum_{1\leq a\leq2s}\sum_{1\leq j,k\leq2s}h_{1ja}n_{a}%
S_{N-1}\left(  a\right)  \delta_{ak}\left\vert \psi_{j}\right\rangle
\left\langle \psi_{k}\right\vert \right\}  \left\vert \psi_{\mu}\right\rangle
\end{align}%
\begin{align}
\frac{\delta}{\delta\psi_{\mu}^{\ast}}\mathcal{H}_{1}\left(  \mathbf{n}%
,\mathbf{\psi},\mathbf{\theta}\right)   &  =\left\{  h_{1}\sum_{1\leq a\leq
2s}n_{a}S_{N-1}\left(  a\right)  \left\vert \psi_{a}\right\rangle \left\langle
\psi_{a}\right\vert \right\}  \left\vert \psi_{\mu}\right\rangle \\
\frac{\delta}{\delta\psi_{\mu}^{\ast}}\mathcal{H}_{1}\left(  \mathbf{n}%
,\mathbf{\psi},\mathbf{\theta}\right)   &  =\sum_{1\leq a\leq2s}n_{a}%
S_{N-1}\left(  a\right)  h_{1}\left\vert \psi_{a}\right\rangle \delta_{a\mu
}=n_{\mu}S_{N-1}\left(  \mu\right)  h_{1}\left\vert \psi_{\mu}\right\rangle
\end{align}%
\begin{equation}
\frac{\delta\mathcal{H}_{1}\left(  \mathbf{n},\mathbf{\psi},\mathbf{\theta
}\right)  }{\delta\psi_{\mu}^{\ast}}=n_{\mu}S_{N-1}\left(  \mu\right)
h_{1}\left\vert \psi_{\mu}\right\rangle .
\end{equation}
The two electron functional derivatives are
\begin{subequations}
\begin{align}
&  \frac{\delta}{\delta\psi_{\mu}^{\ast}}\left(  \left.  \psi_{a}\psi
_{b}\right\vert h_{2}\psi_{c}\psi_{d}\right)  =\left(  \left.  \psi
_{b}\right\vert h_{2}\psi_{d}\right)  \left\vert \psi_{c}\right\rangle
\delta_{a\mu}+\left(  \left.  \psi_{a}\right\vert h_{2}\psi_{c}\right)
\left\vert \psi_{d}\right\rangle \delta_{b\mu}\nonumber\\
&  =\left\{  \overset{}{\left(  \left.  \psi_{b}\right\vert h_{2}\psi
_{d}\right)  \left\vert \psi_{c}\right\rangle \left\langle \psi_{a}\right\vert
+\left(  \left.  \psi_{a}\right\vert h_{2}\psi_{c}\right)  \left\vert \psi
_{d}\right\rangle \left\langle \psi_{b}\right\vert }\right\}  \left\vert
\psi_{\mu}\right\rangle . \label{2partfuncder_o}%
\end{align}
Noting that
\end{subequations}
\begin{equation}
\left(  \left.  \psi_{b}\right\vert h_{2}\psi_{d}\right)  _{kl}\equiv
h_{2kbld}=\int\psi_{k}^{\ast}\left(  \mathbf{x}_{1}\right)  \psi_{b}^{\ast
}\left(  \mathbf{x}_{2}\right)  \frac{1}{\left\Vert \mathbf{r}_{1}%
-\mathbf{r}_{2}\right\Vert }\psi_{l}\left(  \mathbf{x}_{1}\right)  \psi
_{d}\left(  \mathbf{x}_{2}\right)  d\mathbf{x}_{1}d\mathbf{x}_{2},
\end{equation}
we obtain%
\begin{equation}
\left(  \left.  \psi_{b}\right\vert h_{2}\psi_{d}\right)  =\sum_{1\leq
k,l\leq2s}h_{2kbld}\left\vert \psi_{k}\right\rangle \left\langle \psi
_{l}\right\vert =\sum_{1\leq k,l\leq2s}h_{2bkdl}\left\vert \psi_{k}%
\right\rangle \left\langle \psi_{l}\right\vert =\sum_{1\leq k,l\leq
2s}h_{2ldkb}^{\ast}\left\vert \psi_{k}\right\rangle \left\langle \psi
_{l}\right\vert .
\end{equation}
The partial derivatives of $\left\{  \left(  \left.  \psi_{b}\right\vert
h_{2}\psi_{d}\right)  \right\}  $ wrt $\left\{  \psi_{\mu}^{\ast}\right\}  $
are
\begin{align}
&  \frac{\delta}{\delta\psi_{\mu}^{\ast}}\left(  \left.  \psi_{b}\right\vert
h_{2}\psi_{d}\right)  \left.  =\right.  \sum_{1\leq k,l\leq2s}\left\{
h_{2lbmd}\left\vert \psi_{k}\right\rangle \left\langle \left.  \psi
_{l}\right\vert \psi_{c}\right\rangle \left\langle \psi_{a}\right\vert
+h_{2lamc}\left\vert \psi_{k}\right\rangle \left\langle \left.  \psi
_{l}\right\vert \psi_{d}\right\rangle \left\langle \psi_{b}\right\vert
\right\} \nonumber\\
&  \left.  =\right.  \sum_{1\leq k\leq2s}\left\{  h_{2kbcd}\left\vert \psi
_{k}\right\rangle \left\langle \psi_{a}\right\vert +h_{2kadc}\left\vert
\psi_{k}\right\rangle \left\langle \psi_{b}\right\vert \right\}  =\sum_{1\leq
k,l\leq2s}\left\{  h_{2kbcd}\delta_{al}+h_{2kadc}\delta_{bl}\right\}
\left\vert \psi_{k}\right\rangle \left\langle \psi_{l}\right\vert .
\label{2partfuncder_ex}%
\end{align}
The partial functional derivatives of $\mathcal{H}_{2}\left(  \mathbf{\psi
,n}\right)  $ wrt CGSO's $\mathbf{\psi}^{\ast}$ is then given by for $1\leq
\mu\leq s$%
\begin{align}
&  \frac{\delta\mathcal{H}_{2}\left(  \mathbf{n},\mathbf{\psi},\mathbf{\theta
}\right)  }{\delta\psi_{\mu}^{\ast}}=n_{\mu}S_{N-1}\left(  \mu\right)
\left\{  \left(  \left.  \psi_{\mu+s}\right\vert h_{2}\psi_{\mu+s}\right)
\left\vert \psi_{\mu}\right\rangle -\left(  \left.  \psi_{\mu+s}\right\vert
h_{2}\psi_{\mu}\right)  \left\vert \psi_{\mu+s}\right\rangle \right\}
\nonumber\\
&  +\sum_{1\leq k\leq s}n_{\mu}^{\frac{1}{2}}n_{k}^{\frac{1}{2}}e^{i\left(
\theta_{\mu}-\theta_{k}\right)  }S_{N-1}\left(  \mu k\right)  \left\{  \left(
\left.  \psi_{\mu+s}\right\vert h_{2}\psi_{k+s}\right)  \left\vert \psi
_{k}\right\rangle -\left(  \left.  \psi_{\mu+s}\right\vert h_{2}\psi
_{k}\right)  \left\vert \psi_{k+s}\right\rangle \right\} \nonumber\\
&  +\sum_{1\leq k\leq2s}n_{\mu}n_{k}S_{N-2}\left(  \mu k\right)  \left\{
\overset{}{\left(  \left.  \psi_{k}\right\vert h_{2}\psi_{k}\right)
\left\vert \psi_{\mu}\right\rangle }-\left(  \left.  \psi_{k}\right\vert
h_{2}\psi_{\mu}\right)  \left\vert \psi_{k}\right\rangle \right\}  .
\end{align}%
\begin{align}
&  n_{\mu}S_{N-1}\left(  \mu\right)  \left\{  \left(  \left.  \psi_{\mu
+s}\right\vert h_{2}\psi_{\mu+s}\right)  \left\vert \psi_{\mu}\right\rangle
-\left(  \left.  \psi_{\mu+s}\right\vert h_{2}\psi_{\mu}\right)  \left\vert
\psi_{\mu+s}\right\rangle \right\} \\
&  =
\end{align}
for $s+1\leq\mu\leq2s$%
\begin{align}
&  \frac{\delta\mathcal{H}_{2}\left(  \mathbf{n},\mathbf{\psi},\mathbf{\theta
}\right)  }{\delta\psi_{\mu}^{\ast}}=n_{\mu}S_{N-1}\left(  \mu\right)
\left\{  \left(  \left.  \psi_{\mu-s}\right\vert h_{2}\psi_{\mu-s}\right)
\left\vert \psi_{\mu}\right\rangle -\left(  \left.  \psi_{\mu-s}\right\vert
h_{2}\psi_{\mu}\right)  \left\vert \psi_{\mu-s}\right\rangle \right\}
\nonumber\\
&  +\sum_{1\leq k\leq s}n_{\mu}^{\frac{1}{2}}n_{k}^{\frac{1}{2}}e^{i\left(
\theta_{\mu}-\theta_{k}\right)  }S_{N-1}\left(  \mu k\right)  \left\{  \left(
\left.  \psi_{\mu-s}\right\vert h_{2}\psi_{k}\right)  \left\vert \psi
_{k+s}\right\rangle -\left(  \left.  \psi_{\mu-s}\right\vert h_{2}\psi
_{k+s}\right)  \left\vert \psi_{k}\right\rangle \right\} \nonumber\\
&  +\sum_{1\leq k\leq2s}n_{\mu}n_{k}S_{N-2}\left(  \mu k\right)  \left\{
\overset{}{\left(  \left.  \psi_{k}\right\vert h_{2}\psi_{k}\right)
\left\vert \psi_{\mu}\right\rangle }-\left(  \left.  \psi_{k}\right\vert
h_{2}\psi_{\mu}\right)  \left\vert \psi_{k}\right\rangle \right\}  .
\end{align}


We can then define a generalized Fock operator
\begin{align}
&  F\left(  \mathbf{n},\mathbf{\psi},\mathbf{\theta}\right)  =\sum
_{\substack{1\leq j\leq s\\1\leq m\leq2s}}n_{j}S_{N-1}\left(  \jmath\right)
\left\{  \left\vert \psi_{m}\right\rangle \left\langle \psi_{j}\right\vert
h_{1mj}+\left\vert \psi_{j+s}\right\rangle \left\langle \psi_{j+s}\right\vert
h_{1mj+s}+\left\vert \psi_{m}\right\rangle \left\langle \psi_{j}\right\vert
\left(  \left.  \psi_{j+s}\right\vert h_{2}\psi_{j+s}\right)  _{mj}\right.
\nonumber\\
&  \qquad\left.  +\,\left\vert \psi_{m}\right\rangle \left\langle \psi
_{j+s}\right\vert \left(  \left.  \psi_{j}\right\vert h_{2}\psi_{j}\right)
_{mj+s}-\left\vert \psi_{m}\right\rangle \left\langle \psi_{j}\right\vert
\left(  \left.  \psi_{j+s}\right\vert h_{2}\psi_{j}\right)  _{mj+s}-\left\vert
\psi_{m}\right\rangle \left\langle \psi_{j+s}\right\vert \left(  \left.
\psi_{j}\right\vert h_{2}\psi_{j+s}\right)  _{mj}\right\} \nonumber\\
&  \qquad+\sum_{\substack{1\leq j,k\leq s\\1\leq m\leq2s}}n_{j}^{\frac{1}{2}%
}n_{k}^{\frac{1}{2}}e^{i\left(  \theta_{j}-\theta_{k}\right)  }S_{N-1}\left(
\jmath k\right)  \left\{  \left\vert \psi_{m}\right\rangle \left\langle
\psi_{j}\right\vert \left(  \left.  \psi_{j+s}\right\vert h_{2}\psi
_{k+s}\right)  _{mk}+\left\vert \psi_{m}\right\rangle \left\langle \psi
_{j+s}\right\vert \left(  \left.  \psi_{j}\right\vert h_{2}\psi_{k}\right)
_{mk+s}\right. \nonumber\\
&  \qquad\qquad\qquad\qquad\qquad\qquad\left.  -\left\vert \psi_{m}%
\right\rangle \left\langle \psi_{j}\right\vert \left(  \left.  \psi
_{j+s}\right\vert h_{2}\psi_{k}\right)  _{mk+s}-\left\vert \psi_{m}%
\right\rangle \left\langle \psi_{j+s}\right\vert \left(  \left.  \psi
_{j}\right\vert h_{2}\psi_{k+s}\right)  _{mk}\right\} \nonumber\\
&  \qquad+\sum_{\substack{1\leq j,k\leq2s\\1\leq m\leq2s}}n_{j}n_{k}%
S_{N-2}\left(  \jmath k\right)  \left\{  \left\vert \psi_{m}\right\rangle
\left\langle \psi_{j}\right\vert \left(  \left.  \psi_{k}\right\vert h_{2}%
\psi_{k}\right)  _{mj}-\left\vert \psi_{m}\right\rangle \left\langle \psi
_{j}\right\vert \left(  \left.  \psi_{k}\right\vert h_{2}\psi_{j}\right)
_{mk}\right\}  ,
\end{align}
giving%
\begin{equation}
F\left(  \mathbf{n},\mathbf{\psi},\mathbf{\theta}\right)  =\sum_{1\leq
j,m\leq2s}\left\{  n_{j}S_{N-1}\left(  \jmath\right)  \left\{  h_{1mj}%
+A_{mj}-B_{mj}\right\}  +\left\{  C_{mj}-D_{mj}\right\}  +\overset{}{\left\{
E_{mj}-F_{mj}\right\}  \left\vert \psi_{m}\right\rangle \left\langle \psi
_{j}\right\vert }\right\}  .
\end{equation}
The matrix elements of $\left\{  \mathbb{A},\mathbb{B},\mathbb{C}%
,\mathbb{D},\mathbb{E},\mathbb{F}\right\}  $ for $1\leq m\leq2s$ are given by
\begin{equation}
A_{mj}=\left\{
\begin{array}
[c]{c}%
\left(  \left.  \psi_{j+s}\right\vert h_{2}\psi_{j+s}\right)  _{mj};1\leq
j\leq s\\
\left(  \left.  \psi_{j-s}\right\vert h_{2}\psi_{j-s}\right)  _{mj};s+1\leq
j\leq s
\end{array}
\right.  ,
\end{equation}%
\begin{equation}
B_{mj}=\left\{
\begin{array}
[c]{c}%
\left(  \left.  \psi_{j+s}\right\vert h_{2}\psi_{j}\right)  _{mj+s};1\leq
j\leq s\\
\left(  \left.  \psi_{j-s}\right\vert h_{2}\psi_{j}\right)  _{mj-s};s+1\leq
j\leq s
\end{array}
\right.  ,
\end{equation}
noting that $n_{j+s}=n_{j}$ and $\theta_{j+s}=\theta_{j}$ for $1\leq j\leq s$
we have the phase dependent $2$-electron terms
\begin{equation}
C_{mj}=\left\{
\begin{array}
[c]{c}%
%TCIMACRO{\tsum \limits_{1\leq k\leq s}}%
%BeginExpansion
{\textstyle\sum\limits_{1\leq k\leq s}}
%EndExpansion
n_{j}^{\frac{1}{2}}n_{k}^{\frac{1}{2}}S_{N-1}\left(  \jmath k\right)
e^{i\left(  \theta_{j}-\theta_{k}\right)  }\left(  \left.  \psi_{j+s}%
\right\vert h_{2}\psi_{k+s}\right)  _{mk};1\leq j\leq s\\%
%TCIMACRO{\tsum \limits_{1\leq k\leq s}}%
%BeginExpansion
{\textstyle\sum\limits_{1\leq k\leq s}}
%EndExpansion
n_{j}^{\frac{1}{2}}n_{k}^{\frac{1}{2}}S_{N-1}\left(  \jmath k\right)
e^{i\left(  \theta_{j}-\theta_{k}\right)  }\left(  \left.  \psi_{j-s}%
\right\vert h_{2}\psi_{k}\right)  _{mk+s};s+1\leq j\leq2s
\end{array}
\right.  ,
\end{equation}%
\begin{equation}
D_{mj}=\left\{
\begin{array}
[c]{c}%
%TCIMACRO{\tsum \limits_{1\leq k\leq s}}%
%BeginExpansion
{\textstyle\sum\limits_{1\leq k\leq s}}
%EndExpansion
n_{j}^{\frac{1}{2}}n_{k}^{\frac{1}{2}}S_{N-1}\left(  \jmath k\right)
e^{i\left(  \theta_{j}-\theta_{k}\right)  }\left(  \left.  \psi_{j+s}%
\right\vert h_{2}\psi_{k}\right)  _{mk+s};1\leq j\leq s\\%
%TCIMACRO{\tsum \limits_{1\leq k\leq s}}%
%BeginExpansion
{\textstyle\sum\limits_{1\leq k\leq s}}
%EndExpansion
n_{j}^{\frac{1}{2}}n_{k}^{\frac{1}{2}}S_{N-1}\left(  \jmath k\right)
e^{i\left(  \theta_{j}-\theta_{k}\right)  }\left(  \left.  \psi_{j-s}%
\right\vert h_{2}\psi_{k-s}\right)  _{mk};s+1\leq j\leq2s
\end{array}
\right.  ,
\end{equation}
the phase independent $2$-electron terms
\begin{equation}
F_{mj}=\sum_{1\leq k\leq2s}n_{j}n_{k}S_{N-2}\left(  \jmath k\right)  \left(
\left.  \psi_{k}\right\vert h_{2}\psi_{j}\right)  _{mk};1\leq j\leq2s,
\end{equation}
and%
\begin{equation}
E_{mj}=\sum_{1\leq k\leq2s}n_{j}n_{k}S_{N-2}\left(  \jmath k\right)  \left(
\left.  \psi_{k}\right\vert h_{2}\psi_{k}\right)  _{mj};1\leq j\leq2s.
\end{equation}


The Fock operator $F\left(  \mathbf{n},\mathbf{\psi},\mathbf{\theta}\right)  $
can be seen to be self adjoint by considering $\left\langle \psi_{a}\left\vert
F\left(  \mathbf{n},\mathbf{\psi},\mathbf{\theta}\right)  \psi_{b}\right.
\right\rangle $ and $\left\langle \psi_{a}\left\vert F\left(  \mathbf{n}%
,\mathbf{\psi},\mathbf{\theta}\right)  ^{\dagger}\psi_{b}\right.
\right\rangle $ for $1\leq a,b\leq2s,$ where
\begin{align}
&  \left\langle \psi_{a}\left\vert F\left(  \mathbf{n},\mathbf{\psi
},\mathbf{\theta}\right)  \psi_{b}\right.  \right\rangle =\left\langle
\overset{\overset{}{}}{\psi_{a}}\right\vert \left\{  \sum_{1\leq j\leq s}%
n_{j}S_{N-1}\left(  \jmath\right)  \left\{  h_{1}\left\{  \left\vert \psi
_{j}\right\rangle \left\langle \psi_{j}\right\vert +\left\vert \psi
_{j+s}\right\rangle \left\langle \psi_{j+s}\right\vert \right\}  +\left(
\left.  \psi_{j+s}\right\vert h_{2}\psi_{j+s}\right)  \left\vert \psi
_{j}\right\rangle \left\langle \psi_{j}\right\vert \right.  \right.
\nonumber\\
&  \left.  +\,\left(  \left.  \psi_{j}\right\vert h_{2}\psi_{j}\right)
\left\vert \psi_{j+s}\right\rangle \left\langle \psi_{j+s}\right\vert -\left(
\left.  \psi_{j+s}\right\vert h_{2}\psi_{j}\right)  \left\vert \psi
_{j+s}\right\rangle \left\langle \psi_{j}\right\vert -\left(  \left.  \psi
_{j}\right\vert h_{2}\psi_{j+s}\right)  \left\vert \psi_{j}\right\rangle
\left\langle \psi_{j+s}\right\vert \right\} \nonumber\\
&  +\sum_{1\leq j,k\leq s}n_{j}^{\frac{1}{2}}n_{k}^{\frac{1}{2}}e^{i\left(
\theta_{j}-\theta_{k}\right)  }S_{N-1}\left(  \jmath k\right)  \left\{
\left(  \left.  \psi_{j+s}\right\vert h_{2}\psi_{k+s}\right)  \left\vert
\psi_{k}\right\rangle \left\langle \psi_{j}\right\vert +\left(  \left.
\psi_{j}\right\vert h_{2}\psi_{k}\right)  \left\vert \psi_{k+s}\right\rangle
\left\langle \psi_{j+s}\right\vert \right. \nonumber\\
&  \left.  -\left(  \left.  \psi_{j+s}\right\vert h_{2}\psi_{k}\right)
\left\vert \psi_{k+s}\right\rangle \left\langle \psi_{j}\right\vert -\left(
\left.  \psi_{j}\right\vert h_{2}\psi_{k+s}\right)  \left\vert \psi
_{k}\right\rangle \left\langle \psi_{j+s}\right\vert \right\} \nonumber\\
&  \left.  \left.  +\sum_{1\leq j,k\leq2s}n_{j}n_{k}S_{N-2}\left(  \jmath
k\right)  \overset{}{\left\{  \left(  \left.  \psi_{k}\right\vert h_{2}%
\psi_{k}\right)  \left\vert \psi_{j}\right\rangle \left\langle \psi
_{j}\right\vert -\left(  \left.  \psi_{k}\right\vert h_{2}\psi_{j}\right)
\left\vert \psi_{k}\right\rangle \left\langle \psi_{j}\right\vert \right\}
}\right\}  \psi_{b}\right\rangle \label{adj1}%
\end{align}
and
\begin{align}
&  \left\langle \psi_{a}\left\vert F\left(  \mathbf{n},\mathbf{\psi
},\mathbf{\theta}\right)  ^{\dagger}\psi_{b}\right.  \right\rangle
=\left\langle \overset{\overset{}{}}{\psi_{a}}\right\vert \left\{  \sum_{1\leq
j\leq s}n_{j}S_{N-1}\left(  \jmath\right)  \left\{  \left\{  \left\vert
\psi_{j}\right\rangle \left\langle \psi_{j}\right\vert +\left\vert \psi
_{j+s}\right\rangle \left\langle \psi_{j+s}\right\vert \right\}
h_{1}+\left\vert \psi_{j}\right\rangle \left\langle \psi_{j}\right\vert
\left(  \left.  \psi_{j+s}\right\vert h_{2}\psi_{j+s}\right)  \right.  \right.
\nonumber\\
&  \left.  +\left\vert \psi_{j+s}\right\rangle \left\langle \psi
_{j+s}\right\vert \left(  \left.  \psi_{j}\right\vert h_{2}\psi_{j}\right)
-\left\vert \psi_{j}\right\rangle \left\langle \psi_{j+s}\right\vert \left(
\left.  \psi_{j}\right\vert h_{2}\psi_{j+s}\right)  -\left\vert \psi
_{j+s}\right\rangle \left\langle \psi_{j}\right\vert \left(  \left.
\psi_{j+s}\right\vert h_{2}\psi_{j}\right)  \right\} \nonumber\\
&  +\sum_{1\leq j,k\leq s}n_{j}^{\frac{1}{2}}n_{k}^{\frac{1}{2}}e^{i\left(
\theta_{k}-\theta_{j}\right)  }S_{N-1}\left(  \jmath k\right)  \left\{
\left\vert \psi_{j}\right\rangle \left\langle \psi_{k}\right\vert \left(
\left.  \psi_{k+s}\right\vert h_{2}\psi_{j+s}\right)  +\left\vert \psi
_{j+s}\right\rangle \left\langle \psi_{k+s}\right\vert \left(  \left.
\psi_{k}\right\vert h_{2}\psi_{j}\right)  \right. \nonumber\\
&  \left.  -\left\vert \psi_{j}\right\rangle \left\langle \psi_{k+s}%
\right\vert \left(  \left.  \psi_{k}\right\vert h_{2}\psi_{j+s}\right)
-\left\vert \psi_{j+s}\right\rangle \left\langle \psi_{k}\right\vert \left(
\left.  \psi_{k+s}\right\vert h_{2}\psi_{j}\right)  \right\} \nonumber\\
&  \left.  \left.  +\sum_{1\leq j,k\leq2s}n_{j}n_{k}S_{N-2}\left(  \jmath
k\right)  \overset{}{\left\{  \left\vert \psi_{j}\right\rangle \left\langle
\psi_{j}\right\vert \left(  \left.  \psi_{k}\right\vert h_{2}\psi_{k}\right)
-\left\vert \psi_{j}\right\rangle \left\langle \psi_{k}\right\vert \left(
\left.  \psi_{j}\right\vert h_{2}\psi_{k}\right)  \right\}  }\right\}
\psi_{b}\right\rangle . \label{adj2}%
\end{align}
Eq. $\left(  \ref{adj1}\right)  $ gives for $1\leq a\leq2s,1\leq b\leq s$
\begin{align}
&  \left\langle \psi_{a}\left\vert F\left(  \mathbf{n},\mathbf{\psi
},\mathbf{\theta}\right)  \psi_{b}\right.  \right\rangle =n_{b}S_{N-1}\left(
b\right)  \left\{  h_{1ab}+T_{ab+sbb+s}-T_{ab+sb+sb}\right\} \nonumber\\
&  +\sum_{1\leq k\leq s}n_{b}^{\frac{1}{2}}n_{k}^{\frac{1}{2}}e^{i\left(
\theta_{b}-\theta_{k}\right)  }S_{N-1}\left(  bk\right)  \left\{
T_{ab+skk+s}-T_{ab+sk+sk}\right\}  +\sum_{1\leq k\leq2s}n_{b}n_{k}%
S_{N-2}\left(  \jmath k\right)  \left\{  T_{akbk}-T_{akkb}\right\}
\end{align}
and for $1\leq a\leq2s,s+1\leq b\leq2s$%
\begin{align}
&  \left\langle \psi_{a}\left\vert F\left(  \mathbf{n},\mathbf{\psi
},\mathbf{\theta}\right)  \psi_{b}\right.  \right\rangle =n_{b}S_{N-1}\left(
b\right)  \left\{  h_{1ab}+\,T_{ab-sbb-s}-T_{ab-sb-sb}\right\} \nonumber\\
&  +\sum_{1\leq k\leq s}n_{b}^{\frac{1}{2}}n_{k}^{\frac{1}{2}}e^{i\left(
\theta_{b}-\theta_{k}\right)  }S_{N-1}\left(  bk\right)  \left\{
T_{ab-sk+sk}-T_{ab-skk+s}\right\}  +\sum_{1\leq k\leq2s}n_{b}n_{k}%
S_{N-2}\left(  bk\right)  \overset{}{\left\{  T_{akbk}-T_{akkb}\right\}  }.
\end{align}
Eq. $\left(  \ref{adj2}\right)  $ gives for $1\leq a\leq2s,1\leq b\leq s$%
\begin{align}
&  \left\langle \psi_{a}\left\vert F\left(  \mathbf{n},\mathbf{\psi
},\mathbf{\theta}\right)  ^{\dagger}\psi_{b}\right.  \right\rangle
=n_{j}S_{N-1}\left(  \jmath\right)  \left\{  h_{1ab}+T_{aa+sba+s}%
-T_{a+saba+s}\right\} \nonumber\\
&  +\sum_{1\leq k\leq s}n_{j}^{\frac{1}{2}}n_{k}^{\frac{1}{2}}e^{i\left(
\theta_{k}-\theta_{j}\right)  }S_{N-1}\left(  \jmath k\right)  \left\{
T_{kk+sba+s}-T_{k+skba+s}\right\}  +\sum_{1\leq k\leq2s}n_{j}n_{k}%
S_{N-2}\left(  \jmath k\right)  \left\{  T_{akbk}-T_{kabk}\right\}
\end{align}
and for $1\leq a\leq2s,s+1\leq b\leq2s$%
\begin{align}
&  \left\langle \psi_{a}\left\vert F\left(  \mathbf{n},\mathbf{\psi
},\mathbf{\theta}\right)  ^{\dagger}\psi_{b}\right.  \right\rangle
=n_{j}S_{N-1}\left(  \jmath\right)  \left\{  h_{1ab}+T_{aa-sba}-T_{aa-sba}%
\right\} \nonumber\\
&  +\sum_{1\leq j,k\leq s}n_{j}^{\frac{1}{2}}n_{k}^{\frac{1}{2}}e^{i\left(
\theta_{k}-\theta_{j}\right)  }S_{N-1}\left(  \jmath k\right)  \left\{
T_{k+ska-sb}-T_{kk+sba-s}\right\}  +\sum_{1\leq k\leq2s}n_{j}n_{k}%
S_{N-2}\left(  \jmath k\right)  \left\{  T_{akbk}-T_{kabk}\right\}
\end{align}


The derivatives eqs. $\left(  \ref{der1}\right)  $ and eqs. $\left(
\ref{der2}\right)  $ can be expressed as
\begin{align}
\frac{\delta\mathcal{H}\left(  \mathbf{\psi,n}\right)  }{\delta\psi_{\mu
}^{\ast}}  &  =F\left(  \mathbf{n},\mathbf{\psi},\mathbf{\theta}\right)
\left\vert \psi_{\mu}\right\rangle \\
\frac{\delta\mathcal{H}\left(  \mathbf{\psi,n}\right)  }{\delta\psi_{\mu}}  &
=\left\langle \psi_{\mu}\right\vert F\left(  \mathbf{n},\mathbf{\psi
},\mathbf{\theta}\right)  .
\end{align}


\section{Fock Operator and Energy}

Polynomial functions can be expanded as a sum of homogeneous mononomial
functions as%

\begin{equation}
f\left(  \mathbf{x}\right)  =f_{0}+f_{1}\left(  \mathbf{x}\right)
+\cdots+f_{n}\left(  \mathbf{x}\right)  =%
%TCIMACRO{\tsum \limits_{0\leq m\leq n}}%
%BeginExpansion
{\textstyle\sum\limits_{0\leq m\leq n}}
%EndExpansion
f_{m}\left(  \mathbf{x}\right)  ,
\end{equation}
where
\begin{equation}
f_{m}\left(  \mathbf{x}\right)  =%
%TCIMACRO{\tsum \limits_{1\leq j_{1},\cdots,j_{m}\leq r}}%
%BeginExpansion
{\textstyle\sum\limits_{1\leq j_{1},\cdots,j_{m}\leq r}}
%EndExpansion
c_{j_{1}\cdots j_{m}}x_{j_{1}}\cdots x_{j_{m}}.
\end{equation}
Consider
\begin{align}
\frac{\partial f_{m}}{\partial x_{l}}\left(  \mathbf{x}\right)   &  =%
%TCIMACRO{\tsum \limits_{1\leq j_{1},\cdots,j_{m-1}\neq l\leq r}}%
%BeginExpansion
{\textstyle\sum\limits_{1\leq j_{1},\cdots,j_{m-1}\neq l\leq r}}
%EndExpansion
c_{j_{1}\cdots j_{l}\cdots j_{m}}x_{j_{1}}\cdots x_{j_{m-1}}\nonumber\\
x_{l}\frac{\partial f_{m}}{\partial x_{l}}\left(  \mathbf{x}\right)   &  =%
%TCIMACRO{\tsum \limits_{1\leq j_{1},\cdots,j_{m-1}\neq l\leq r}}%
%BeginExpansion
{\textstyle\sum\limits_{1\leq j_{1},\cdots,j_{m-1}\neq l\leq r}}
%EndExpansion
c_{j_{1}\cdots j_{l}\cdots j_{m}}x_{j_{1}}\cdots x_{j_{m-1}}x_{l}\nonumber\\%
%TCIMACRO{\tsum \limits_{1\leq l\leq r}}%
%BeginExpansion
{\textstyle\sum\limits_{1\leq l\leq r}}
%EndExpansion
x_{l}\frac{\partial f_{m}}{\partial x_{l}}\left(  \mathbf{x}\right)   &
=mf_{m}\left(  \mathbf{x}\right)
\end{align}
hence%
\begin{align}%
%TCIMACRO{\tsum \limits_{1\leq l\leq r}}%
%BeginExpansion
{\textstyle\sum\limits_{1\leq l\leq r}}
%EndExpansion
x_{l}\frac{\partial f}{\partial x_{l}}\left(  \mathbf{x}\right)   &  =%
%TCIMACRO{\tsum \limits_{1\leq l\leq r}}%
%BeginExpansion
{\textstyle\sum\limits_{1\leq l\leq r}}
%EndExpansion
x_{l}\frac{\partial}{\partial x_{l}}\left\{
%TCIMACRO{\tsum \limits_{0\leq m\leq n}}%
%BeginExpansion
{\textstyle\sum\limits_{0\leq m\leq n}}
%EndExpansion
f_{m}\left(  \mathbf{x}\right)  \right\}  =%
%TCIMACRO{\tsum \limits_{\substack{1\leq l\leq r\\0\leq m\leq n}}}%
%BeginExpansion
{\textstyle\sum\limits_{\substack{1\leq l\leq r\\0\leq m\leq n}}}
%EndExpansion
x_{l}\frac{\partial f_{m}\left(  \mathbf{x}\right)  }{\partial x_{l}%
}\nonumber\\
&  =%
%TCIMACRO{\tsum \limits_{\substack{0\leq m\leq n\\1\leq l\leq r}}}%
%BeginExpansion
{\textstyle\sum\limits_{\substack{0\leq m\leq n\\1\leq l\leq r}}}
%EndExpansion
x_{l}\frac{\partial f_{m}\left(  \mathbf{x}\right)  }{\partial x_{l}}=%
%TCIMACRO{\tsum \limits_{0\leq m\leq n}}%
%BeginExpansion
{\textstyle\sum\limits_{0\leq m\leq n}}
%EndExpansion
mf_{m}\left(  \mathbf{x}\right) \nonumber\\
&  =f_{1}\left(  \mathbf{x}\right)  +2f_{2}\left(  \mathbf{x}\right)
+\cdots+nf_{n}\left(  \mathbf{x}\right)  .
\end{align}
Thus
\begin{equation}%
%TCIMACRO{\tsum \limits_{1\leq l\leq2s}}%
%BeginExpansion
{\textstyle\sum\limits_{1\leq l\leq2s}}
%EndExpansion
\left\langle \left.  \psi_{l}\right\vert F\left(  \mathbf{\psi,n}\right)
\psi_{l}\right\rangle =Tr\left\{  F\left(  \mathbf{\psi,n}\right)  \right\}
=\mathcal{H}_{1}\left(  \mathbf{\psi,n}\right)  +2\mathcal{H}_{2}\left(
\mathbf{\psi,n}\right)
\end{equation}
showing that
\begin{equation}
Tr\left\{  F\left(  \mathbf{\psi,n}\right)  \right\}  -\mathcal{H}\left(
\mathbf{\psi,n}\right)  =\mathcal{H}_{2}\left(  \mathbf{\psi,n}\right)
\end{equation}
and
\begin{equation}
2\mathcal{H}\left(  \mathbf{\psi,n}\right)  =Tr\left\{  F\left(
\mathbf{\psi,n}\right)  \right\}  +\mathcal{H}_{1}\left(  \mathbf{\psi
,n}\right)
\end{equation}
which gives
\begin{equation}
\mathcal{H}\left(  \mathbf{\psi,n}\right)  =Tr\left\{  F\left(  \mathbf{\psi
,n}\right)  \right\}  -\mathcal{H}_{2}\left(  \mathbf{\psi,n}\right)
\end{equation}
and
\begin{equation}
\mathcal{H}\left(  \mathbf{\psi,n}\right)  =\frac{1}{2}\left\{  Tr\left\{
F\left(  \mathbf{\psi,n}\right)  \right\}  +\mathcal{H}_{1}\left(
\mathbf{\psi,n}\right)  \right\}  .
\end{equation}
Noting that
\begin{equation}
Tr\left\{  F\left(  \mathbf{\psi,n}\right)  \right\}  =%
%TCIMACRO{\tsum \limits_{1\leq l\leq2s}}%
%BeginExpansion
{\textstyle\sum\limits_{1\leq l\leq2s}}
%EndExpansion
\varepsilon_{l}\left(  \mathbf{\psi,n}\right)
\end{equation}
and
\begin{equation}
\mathcal{H}_{1}\left(  \mathbf{\psi,n}\right)  =%
%TCIMACRO{\tsum \limits_{1\leq l\leq2s}}%
%BeginExpansion
{\textstyle\sum\limits_{1\leq l\leq2s}}
%EndExpansion
N_{l}\left(  \left.  \psi_{l}\right\vert h_{1}\psi_{l}\right)
\end{equation}
and
\begin{equation}
\left(  \left.  \psi_{l}\right\vert h_{1}\psi_{l}\right)  \leq0,
\end{equation}
it seems that the CGSO's should be occupied by an aufbau principle based on
$\left\{  \left(  \left.  \psi_{l}\right\vert h_{1}\psi_{l}\right)  \right\}
,$ i.e.
\begin{equation}
N_{1}\geq N_{2}\geq\cdots\geq N_{s},
\end{equation}
corresponding to
\begin{equation}
\left(  \left.  \psi_{1}\right\vert h_{1}\psi_{1}\right)  \leq\left(  \left.
\psi_{1+s}\right\vert h_{1}\psi_{1+s}\right)  \leq\left(  \left.  \psi
_{2}\right\vert h_{1}\psi_{2}\right)  \leq\left(  \left.  \psi_{2+s}%
\right\vert h_{1}\psi_{2+s}\right)  \leq\cdots\leq\left(  \left.  \psi
_{s}\right\vert h_{1}\psi_{s}\right)  \leq\left(  \left.  \psi_{2s}\right\vert
h_{1}\psi_{2s}\right)  .
\end{equation}



\end{document}